\documentclass{report}

%%%%%%%%%%%%%%%%%%%%%%%%%%%%%%%%%
% PACKAGE IMPORTS
%%%%%%%%%%%%%%%%%%%%%%%%%%%%%%%%%


\usepackage[tmargin=2cm,rmargin=1in,lmargin=1in,margin=0.85in,bmargin=2cm,footskip=.2in]{geometry}
\usepackage{amsmath,amsfonts,amsthm,amssymb,mathtools}
\usepackage[varbb]{newpxmath}
\usepackage{xfrac}
\usepackage[makeroom]{cancel}
\usepackage{mathtools}
\usepackage{bookmark}
\usepackage{enumitem}
\usepackage{hyperref,theoremref}
\hypersetup{
	pdftitle={Assignment},
	colorlinks=true, linkcolor=doc!90,
	bookmarksnumbered=true,
	bookmarksopen=true
}
\usepackage[most,many,breakable]{tcolorbox}
\usepackage{xcolor}
\usepackage{varwidth}
\usepackage{varwidth}
\usepackage{etoolbox}
%\usepackage{authblk}
\usepackage{nameref}
\usepackage{multicol,array}
\usepackage{tikz-cd}
\usepackage[ruled,vlined,linesnumbered]{algorithm2e}
\usepackage{comment} % enables the use of multi-line comments (\ifx \fi) 
\usepackage{import}
\usepackage{xifthen}
\usepackage{pdfpages}
\usepackage{transparent}
\usepackage[utf8]{inputenc}
\usepackage{fontawesome}


\newcommand\mycommfont[1]{\footnotesize\ttfamily\textcolor{blue}{#1}}
\SetCommentSty{mycommfont}
\newcommand{\incfig}[1]{%
    \def\svgwidth{\columnwidth}
    \import{./figures/}{#1.pdf_tex}
}

\usepackage{tikzsymbols}
\renewcommand{\qedsymbol}{\faWhatsapp}


%\usepackage{import}
%\usepackage{xifthen}
%\usepackage{pdfpages}
%\usepackage{transparent}


%%%%%%%%%%%%%%%%%%%%%%%%%%%%%%
% SELF MADE COLORS
%%%%%%%%%%%%%%%%%%%%%%%%%%%%%%



\definecolor{myg}{RGB}{56, 140, 70}
\definecolor{myb}{RGB}{45, 111, 177}
\definecolor{myr}{RGB}{199, 68, 64}
\definecolor{mytheorembg}{HTML}{F2F2F9}
\definecolor{mytheoremfr}{HTML}{00007B}
\definecolor{mylenmabg}{HTML}{FFFAF8}
\definecolor{mylenmafr}{HTML}{983b0f}
\definecolor{mypropbg}{HTML}{f2fbfc}
\definecolor{mypropfr}{HTML}{191971}
\definecolor{myexamplebg}{HTML}{F2FBF8}
\definecolor{myexamplefr}{HTML}{88D6D1}
\definecolor{myexampleti}{HTML}{2A7F7F}
\definecolor{mydefinitbg}{HTML}{E5E5FF}
\definecolor{mydefinitfr}{HTML}{3F3FA3}
\definecolor{notesgreen}{RGB}{0,162,0}
\definecolor{myp}{RGB}{197, 92, 212}
\definecolor{mygr}{HTML}{2C3338}
\definecolor{myred}{RGB}{127,0,0}
\definecolor{myyellow}{RGB}{169,121,69}
\definecolor{myexercisebg}{HTML}{F2FBF8}
\definecolor{myexercisefg}{HTML}{88D6D1}


%%%%%%%%%%%%%%%%%%%%%%%%%%%%
% TCOLORBOX SETUPS
%%%%%%%%%%%%%%%%%%%%%%%%%%%%

\setlength{\parindent}{1cm}
%================================
% THEOREM BOX
%================================

\tcbuselibrary{theorems,skins,hooks}
\newtcbtheorem[number within=section]{Theorem}{Theorem}
{%
	enhanced,
	breakable,
	colback = mytheorembg,
	frame hidden,
	boxrule = 0sp,
	borderline west = {2pt}{0pt}{mytheoremfr},
	sharp corners,
	detach title,
	before upper = \tcbtitle\par\smallskip,
	coltitle = mytheoremfr,
	fonttitle = \bfseries\sffamily,
	description font = \mdseries,
	separator sign none,
	segmentation style={solid, mytheoremfr},
}
{th}

\tcbuselibrary{theorems,skins,hooks}
\newtcbtheorem[number within=chapter]{theorem}{Theorem}
{%
	enhanced,
	breakable,
	colback = mytheorembg,
	frame hidden,
	boxrule = 0sp,
	borderline west = {2pt}{0pt}{mytheoremfr},
	sharp corners,
	detach title,
	before upper = \tcbtitle\par\smallskip,
	coltitle = mytheoremfr,
	fonttitle = \bfseries\sffamily,
	description font = \mdseries,
	separator sign none,
	segmentation style={solid, mytheoremfr},
}
{th}


\tcbuselibrary{theorems,skins,hooks}
\newtcolorbox{Theoremcon}
{%
	enhanced
	,breakable
	,colback = mytheorembg
	,frame hidden
	,boxrule = 0sp
	,borderline west = {2pt}{0pt}{mytheoremfr}
	,sharp corners
	,description font = \mdseries
	,separator sign none
}

%================================
% Corollery
%================================
\tcbuselibrary{theorems,skins,hooks}
\newtcbtheorem[number within=section]{Corollary}{Corollary}
{%
	enhanced
	,breakable
	,colback = myp!10
	,frame hidden
	,boxrule = 0sp
	,borderline west = {2pt}{0pt}{myp!85!black}
	,sharp corners
	,detach title
	,before upper = \tcbtitle\par\smallskip
	,coltitle = myp!85!black
	,fonttitle = \bfseries\sffamily
	,description font = \mdseries
	,separator sign none
	,segmentation style={solid, myp!85!black}
}
{th}
\tcbuselibrary{theorems,skins,hooks}
\newtcbtheorem[number within=chapter]{corollary}{Corollary}
{%
	enhanced
	,breakable
	,colback = myp!10
	,frame hidden
	,boxrule = 0sp
	,borderline west = {2pt}{0pt}{myp!85!black}
	,sharp corners
	,detach title
	,before upper = \tcbtitle\par\smallskip
	,coltitle = myp!85!black
	,fonttitle = \bfseries\sffamily
	,description font = \mdseries
	,separator sign none
	,segmentation style={solid, myp!85!black}
}
{th}


%================================
% LENMA
%================================

\tcbuselibrary{theorems,skins,hooks}
\newtcbtheorem[number within=section]{Lenma}{Lenma}
{%
	enhanced,
	breakable,
	colback = mylenmabg,
	frame hidden,
	boxrule = 0sp,
	borderline west = {2pt}{0pt}{mylenmafr},
	sharp corners,
	detach title,
	before upper = \tcbtitle\par\smallskip,
	coltitle = mylenmafr,
	fonttitle = \bfseries\sffamily,
	description font = \mdseries,
	separator sign none,
	segmentation style={solid, mylenmafr},
}
{th}

\tcbuselibrary{theorems,skins,hooks}
\newtcbtheorem[number within=chapter]{lenma}{Lenma}
{%
	enhanced,
	breakable,
	colback = mylenmabg,
	frame hidden,
	boxrule = 0sp,
	borderline west = {2pt}{0pt}{mylenmafr},
	sharp corners,
	detach title,
	before upper = \tcbtitle\par\smallskip,
	coltitle = mylenmafr,
	fonttitle = \bfseries\sffamily,
	description font = \mdseries,
	separator sign none,
	segmentation style={solid, mylenmafr},
}
{th}


%================================
% PROPOSITION
%================================

\tcbuselibrary{theorems,skins,hooks}
\newtcbtheorem[number within=section]{Prop}{Proposition}
{%
	enhanced,
	breakable,
	colback = mypropbg,
	frame hidden,
	boxrule = 0sp,
	borderline west = {2pt}{0pt}{mypropfr},
	sharp corners,
	detach title,
	before upper = \tcbtitle\par\smallskip,
	coltitle = mypropfr,
	fonttitle = \bfseries\sffamily,
	description font = \mdseries,
	separator sign none,
	segmentation style={solid, mypropfr},
}
{th}

\tcbuselibrary{theorems,skins,hooks}
\newtcbtheorem[number within=chapter]{prop}{Proposition}
{%
	enhanced,
	breakable,
	colback = mypropbg,
	frame hidden,
	boxrule = 0sp,
	borderline west = {2pt}{0pt}{mypropfr},
	sharp corners,
	detach title,
	before upper = \tcbtitle\par\smallskip,
	coltitle = mypropfr,
	fonttitle = \bfseries\sffamily,
	description font = \mdseries,
	separator sign none,
	segmentation style={solid, mypropfr},
}
{th}


%================================
% CLAIM
%================================

\tcbuselibrary{theorems,skins,hooks}
\newtcbtheorem[number within=section]{claim}{Claim}
{%
	enhanced
	,breakable
	,colback = myg!10
	,frame hidden
	,boxrule = 0sp
	,borderline west = {2pt}{0pt}{myg}
	,sharp corners
	,detach title
	,before upper = \tcbtitle\par\smallskip
	,coltitle = myg!85!black
	,fonttitle = \bfseries\sffamily
	,description font = \mdseries
	,separator sign none
	,segmentation style={solid, myg!85!black}
}
{th}



%================================
% Exercise
%================================

\tcbuselibrary{theorems,skins,hooks}
\newtcbtheorem[number within=section]{Exercise}{Exercise}
{%
	enhanced,
	breakable,
	colback = myexercisebg,
	frame hidden,
	boxrule = 0sp,
	borderline west = {2pt}{0pt}{myexercisefg},
	sharp corners,
	detach title,
	before upper = \tcbtitle\par\smallskip,
	coltitle = myexercisefg,
	fonttitle = \bfseries\sffamily,
	description font = \mdseries,
	separator sign none,
	segmentation style={solid, myexercisefg},
}
{th}

\tcbuselibrary{theorems,skins,hooks}
\newtcbtheorem[number within=chapter]{exercise}{Exercise}
{%
	enhanced,
	breakable,
	colback = myexercisebg,
	frame hidden,
	boxrule = 0sp,
	borderline west = {2pt}{0pt}{myexercisefg},
	sharp corners,
	detach title,
	before upper = \tcbtitle\par\smallskip,
	coltitle = myexercisefg,
	fonttitle = \bfseries\sffamily,
	description font = \mdseries,
	separator sign none,
	segmentation style={solid, myexercisefg},
}
{th}

%================================
% EXAMPLE BOX
%================================

\newtcbtheorem[number within=section]{Example}{Example}
{%
	colback = myexamplebg
	,breakable
	,colframe = myexamplefr
	,coltitle = myexampleti
	,boxrule = 1pt
	,sharp corners
	,detach title
	,before upper=\tcbtitle\par\smallskip
	,fonttitle = \bfseries
	,description font = \mdseries
	,separator sign none
	,description delimiters parenthesis
}
{ex}

\newtcbtheorem[number within=chapter]{example}{Example}
{%
	colback = myexamplebg
	,breakable
	,colframe = myexamplefr
	,coltitle = myexampleti
	,boxrule = 1pt
	,sharp corners
	,detach title
	,before upper=\tcbtitle\par\smallskip
	,fonttitle = \bfseries
	,description font = \mdseries
	,separator sign none
	,description delimiters parenthesis
}
{ex}

%================================
% DEFINITION BOX
%================================

\newtcbtheorem[number within=section]{Definition}{Definition}{enhanced,
	before skip=2mm,after skip=2mm, colback=red!5,colframe=red!80!black,boxrule=0.5mm,
	attach boxed title to top left={xshift=1cm,yshift*=1mm-\tcboxedtitleheight}, varwidth boxed title*=-3cm,
	boxed title style={frame code={
					\path[fill=tcbcolback]
					([yshift=-1mm,xshift=-1mm]frame.north west)
					arc[start angle=0,end angle=180,radius=1mm]
					([yshift=-1mm,xshift=1mm]frame.north east)
					arc[start angle=180,end angle=0,radius=1mm];
					\path[left color=tcbcolback!60!black,right color=tcbcolback!60!black,
						middle color=tcbcolback!80!black]
					([xshift=-2mm]frame.north west) -- ([xshift=2mm]frame.north east)
					[rounded corners=1mm]-- ([xshift=1mm,yshift=-1mm]frame.north east)
					-- (frame.south east) -- (frame.south west)
					-- ([xshift=-1mm,yshift=-1mm]frame.north west)
					[sharp corners]-- cycle;
				},interior engine=empty,
		},
	fonttitle=\bfseries,
	title={#2},#1}{def}
\newtcbtheorem[number within=chapter]{definition}{Definition}{enhanced,
	before skip=2mm,after skip=2mm, colback=red!5,colframe=red!80!black,boxrule=0.5mm,
	attach boxed title to top left={xshift=1cm,yshift*=1mm-\tcboxedtitleheight}, varwidth boxed title*=-3cm,
	boxed title style={frame code={
					\path[fill=tcbcolback]
					([yshift=-1mm,xshift=-1mm]frame.north west)
					arc[start angle=0,end angle=180,radius=1mm]
					([yshift=-1mm,xshift=1mm]frame.north east)
					arc[start angle=180,end angle=0,radius=1mm];
					\path[left color=tcbcolback!60!black,right color=tcbcolback!60!black,
						middle color=tcbcolback!80!black]
					([xshift=-2mm]frame.north west) -- ([xshift=2mm]frame.north east)
					[rounded corners=1mm]-- ([xshift=1mm,yshift=-1mm]frame.north east)
					-- (frame.south east) -- (frame.south west)
					-- ([xshift=-1mm,yshift=-1mm]frame.north west)
					[sharp corners]-- cycle;
				},interior engine=empty,
		},
	fonttitle=\bfseries,
	title={#2},#1}{def}



%================================
% Solution BOX
%================================

\makeatletter
\newtcbtheorem{question}{Question}{enhanced,
	breakable,
	colback=white,
	colframe=myb!80!black,
	attach boxed title to top left={yshift*=-\tcboxedtitleheight},
	fonttitle=\bfseries,
	title={#2},
	boxed title size=title,
	boxed title style={%
			sharp corners,
			rounded corners=northwest,
			colback=tcbcolframe,
			boxrule=0pt,
		},
	underlay boxed title={%
			\path[fill=tcbcolframe] (title.south west)--(title.south east)
			to[out=0, in=180] ([xshift=5mm]title.east)--
			(title.center-|frame.east)
			[rounded corners=\kvtcb@arc] |-
			(frame.north) -| cycle;
		},
	#1
}{def}
\makeatother

%================================
% SOLUTION BOX
%================================

\makeatletter
\newtcolorbox{solution}{enhanced,
	breakable,
	colback=white,
	colframe=myg!80!black,
	attach boxed title to top left={yshift*=-\tcboxedtitleheight},
	title=Solution,
	boxed title size=title,
	boxed title style={%
			sharp corners,
			rounded corners=northwest,
			colback=tcbcolframe,
			boxrule=0pt,
		},
	underlay boxed title={%
			\path[fill=tcbcolframe] (title.south west)--(title.south east)
			to[out=0, in=180] ([xshift=5mm]title.east)--
			(title.center-|frame.east)
			[rounded corners=\kvtcb@arc] |-
			(frame.north) -| cycle;
		},
}
\makeatother

%================================
% Question BOX
%================================

\makeatletter
\newtcbtheorem{qstion}{Question}{enhanced,
	breakable,
	colback=white,
	colframe=mygr,
	attach boxed title to top left={yshift*=-\tcboxedtitleheight},
	fonttitle=\bfseries,
	title={#2},
	boxed title size=title,
	boxed title style={%
			sharp corners,
			rounded corners=northwest,
			colback=tcbcolframe,
			boxrule=0pt,
		},
	underlay boxed title={%
			\path[fill=tcbcolframe] (title.south west)--(title.south east)
			to[out=0, in=180] ([xshift=5mm]title.east)--
			(title.center-|frame.east)
			[rounded corners=\kvtcb@arc] |-
			(frame.north) -| cycle;
		},
	#1
}{def}
\makeatother

\newtcbtheorem[number within=chapter]{wconc}{Wrong Concept}{
	breakable,
	enhanced,
	colback=white,
	colframe=myr,
	arc=0pt,
	outer arc=0pt,
	fonttitle=\bfseries\sffamily\large,
	colbacktitle=myr,
	attach boxed title to top left={},
	boxed title style={
			enhanced,
			skin=enhancedfirst jigsaw,
			arc=3pt,
			bottom=0pt,
			interior style={fill=myr}
		},
	#1
}{def}



%================================
% NOTE BOX
%================================

\usetikzlibrary{arrows,calc,shadows.blur}
\tcbuselibrary{skins}
\newtcolorbox{note}[1][]{%
	enhanced jigsaw,
	colback=gray!20!white,%
	colframe=gray!80!black,
	size=small,
	boxrule=1pt,
	title=\textbf{Note:},
	halign title=flush center,
	coltitle=black,
	breakable,
	drop shadow=black!50!white,
	attach boxed title to top left={xshift=1cm,yshift=-\tcboxedtitleheight/2,yshifttext=-\tcboxedtitleheight/2},
	minipage boxed title=1.5cm,
	boxed title style={%
			colback=white,
			size=fbox,
			boxrule=1pt,
			boxsep=2pt,
			underlay={%
					\coordinate (dotA) at ($(interior.west) + (-0.5pt,0)$);
					\coordinate (dotB) at ($(interior.east) + (0.5pt,0)$);
					\begin{scope}
						\clip (interior.north west) rectangle ([xshift=3ex]interior.east);
						\filldraw [white, blur shadow={shadow opacity=60, shadow yshift=-.75ex}, rounded corners=2pt] (interior.north west) rectangle (interior.south east);
					\end{scope}
					\begin{scope}[gray!80!black]
						\fill (dotA) circle (2pt);
						\fill (dotB) circle (2pt);
					\end{scope}
				},
		},
	#1,
}

%%%%%%%%%%%%%%%%%%%%%%%%%%%%%%
% SELF MADE COMMANDS
%%%%%%%%%%%%%%%%%%%%%%%%%%%%%%


\newcommand{\thm}[2]{\begin{Theorem}{#1}{}#2\end{Theorem}}
\newcommand{\cor}[2]{\begin{Corollary}{#1}{}#2\end{Corollary}}
\newcommand{\mlenma}[2]{\begin{Lenma}{#1}{}#2\end{Lenma}}
\newcommand{\mprop}[2]{\begin{Prop}{#1}{}#2\end{Prop}}
\newcommand{\clm}[3]{\begin{claim}{#1}{#2}#3\end{claim}}
\newcommand{\wc}[2]{\begin{wconc}{#1}{}\setlength{\parindent}{1cm}#2\end{wconc}}
\newcommand{\thmcon}[1]{\begin{Theoremcon}{#1}\end{Theoremcon}}
\newcommand{\ex}[2]{\begin{Example}{#1}{}#2\end{Example}}
\newcommand{\dfn}[2]{\begin{Definition}[colbacktitle=red!75!black]{#1}{}#2\end{Definition}}
\newcommand{\dfnc}[2]{\begin{definition}[colbacktitle=red!75!black]{#1}{}#2\end{definition}}
\newcommand{\qs}[2]{\begin{question}{#1}{}#2\end{question}}
\newcommand{\pf}[2]{\begin{myproof}[#1]#2\end{myproof}}
\newcommand{\nt}[1]{\begin{note}#1\end{note}}

\newcommand*\circled[1]{\tikz[baseline=(char.base)]{
		\node[shape=circle,draw,inner sep=1pt] (char) {#1};}}
\newcommand\getcurrentref[1]{%
	\ifnumequal{\value{#1}}{0}
	{??}
	{\the\value{#1}}%
}
\newcommand{\getCurrentSectionNumber}{\getcurrentref{section}}
\newenvironment{myproof}[1][\proofname]{%
	\proof[\bfseries #1: ]%
}{\endproof}

\newcommand{\mclm}[2]{\begin{myclaim}[#1]#2\end{myclaim}}
\newenvironment{myclaim}[1][\claimname]{\proof[\bfseries #1: ]}{}

\newcounter{mylabelcounter}

\makeatletter
\newcommand{\setword}[2]{%
	\phantomsection
	#1\def\@currentlabel{\unexpanded{#1}}\label{#2}%
}
\makeatother




\tikzset{
	symbol/.style={
			draw=none,
			every to/.append style={
					edge node={node [sloped, allow upside down, auto=false]{$#1$}}}
		}
}


% deliminators
\DeclarePairedDelimiter{\abs}{\lvert}{\rvert}
\DeclarePairedDelimiter{\norm}{\lVert}{\rVert}

\DeclarePairedDelimiter{\ceil}{\lceil}{\rceil}
\DeclarePairedDelimiter{\floor}{\lfloor}{\rfloor}
\DeclarePairedDelimiter{\round}{\lfloor}{\rceil}

\newsavebox\diffdbox
\newcommand{\slantedromand}{{\mathpalette\makesl{d}}}
\newcommand{\makesl}[2]{%
\begingroup
\sbox{\diffdbox}{$\mathsurround=0pt#1\mathrm{#2}$}%
\pdfsave
\pdfsetmatrix{1 0 0.2 1}%
\rlap{\usebox{\diffdbox}}%
\pdfrestore
\hskip\wd\diffdbox
\endgroup
}
\newcommand{\dd}[1][]{\ensuremath{\mathop{}\!\ifstrempty{#1}{%
\slantedromand\@ifnextchar^{\hspace{0.2ex}}{\hspace{0.1ex}}}%
{\slantedromand\hspace{0.2ex}^{#1}}}}
\ProvideDocumentCommand\dv{o m g}{%
  \ensuremath{%
    \IfValueTF{#3}{%
      \IfNoValueTF{#1}{%
        \frac{\dd #2}{\dd #3}%
      }{%
        \frac{\dd^{#1} #2}{\dd #3^{#1}}%
      }%
    }{%
      \IfNoValueTF{#1}{%
        \frac{\dd}{\dd #2}%
      }{%
        \frac{\dd^{#1}}{\dd #2^{#1}}%
      }%
    }%
  }%
}
\providecommand*{\pdv}[3][]{\frac{\partial^{#1}#2}{\partial#3^{#1}}}
%  - others
\DeclareMathOperator{\Lap}{\mathcal{L}}
\DeclareMathOperator{\Var}{Var} % varience
\DeclareMathOperator{\Cov}{Cov} % covarience
\DeclareMathOperator{\E}{E} % expected

% Since the amsthm package isn't loaded

% I prefer the slanted \leq
\let\oldleq\leq % save them in case they're every wanted
\let\oldgeq\geq
\renewcommand{\leq}{\leqslant}
\renewcommand{\geq}{\geqslant}

% % redefine matrix env to allow for alignment, use r as default
% \renewcommand*\env@matrix[1][r]{\hskip -\arraycolsep
%     \let\@ifnextchar\new@ifnextchar
%     \array{*\c@MaxMatrixCols #1}}


%\usepackage{framed}
%\usepackage{titletoc}
%\usepackage{etoolbox}
%\usepackage{lmodern}


%\patchcmd{\tableofcontents}{\contentsname}{\sffamily\contentsname}{}{}

%\renewenvironment{leftbar}
%{\def\FrameCommand{\hspace{6em}%
%		{\color{myyellow}\vrule width 2pt depth 6pt}\hspace{1em}}%
%	\MakeFramed{\parshape 1 0cm \dimexpr\textwidth-6em\relax\FrameRestore}\vskip2pt%
%}
%{\endMakeFramed}

%\titlecontents{chapter}
%[0em]{\vspace*{2\baselineskip}}
%{\parbox{4.5em}{%
%		\hfill\Huge\sffamily\bfseries\color{myred}\thecontentspage}%
%	\vspace*{-2.3\baselineskip}\leftbar\textsc{\small\chaptername~\thecontentslabel}\\\sffamily}
%{}{\endleftbar}
%\titlecontents{section}
%[8.4em]
%{\sffamily\contentslabel{3em}}{}{}
%{\hspace{0.5em}\nobreak\itshape\color{myred}\contentspage}
%\titlecontents{subsection}
%[8.4em]
%{\sffamily\contentslabel{3em}}{}{}  
%{\hspace{0.5em}\nobreak\itshape\color{myred}\contentspage}



%%%%%%%%%%%%%%%%%%%%%%%%%%%%%%%%%%%%%%%%%%%
% TABLE OF CONTENTS
%%%%%%%%%%%%%%%%%%%%%%%%%%%%%%%%%%%%%%%%%%%

\usepackage{tikz}
\definecolor{doc}{RGB}{0,60,110}
\usepackage{titletoc}
\contentsmargin{0cm}
\titlecontents{chapter}[3.7pc]
{\addvspace{30pt}%
	\begin{tikzpicture}[remember picture, overlay]%
		\draw[fill=doc!60,draw=doc!60] (-7,-.1) rectangle (-0.9,.5);%
		\pgftext[left,x=-3.5cm,y=0.2cm]{\color{white}\Large\sc\bfseries Chapter\ \thecontentslabel};%
	\end{tikzpicture}\color{doc!60}\large\sc\bfseries}%
{}
{}
{\;\titlerule\;\large\sc\bfseries Page \thecontentspage
	\begin{tikzpicture}[remember picture, overlay]
		\draw[fill=doc!60,draw=doc!60] (2pt,0) rectangle (4,0.1pt);
	\end{tikzpicture}}%
\titlecontents{section}[3.7pc]
{\addvspace{2pt}}
{\contentslabel[\thecontentslabel]{2pc}}
{}
{\hfill\small \thecontentspage}
[]
\titlecontents*{subsection}[3.7pc]
{\addvspace{-1pt}\small}
{}
{}
{\ --- \small\thecontentspage}
[ \textbullet\ ][]

\makeatletter
\renewcommand{\tableofcontents}{%
	\chapter*{%
	  \vspace*{-20\p@}%
	  \begin{tikzpicture}[remember picture, overlay]%
		  \pgftext[right,x=15cm,y=0.2cm]{\color{doc!60}\Huge\sc\bfseries \contentsname};%
		  \draw[fill=doc!60,draw=doc!60] (13,-.75) rectangle (20,1);%
		  \clip (13,-.75) rectangle (20,1);
		  \pgftext[right,x=15cm,y=0.2cm]{\color{white}\Huge\sc\bfseries \contentsname};%
	  \end{tikzpicture}}%
	\@starttoc{toc}}
\makeatother


% Comandi di Bastiality %
\newcommand{\innerproduct}[2]{\langle #1, #2 \rangle}
\newcommand{\generators}[1]{\langle #1 \rangle}
\newcommand{\powerset}[0]{\mathcal{P}}
\newcommand{\indip}{\perp \!\!\! \perp}

% Comandi di GioLaPalma %
\newcommand{\der}[0]{\implies^*}
\newcommand{\red}[1]{\textcolor{red}{#1}}
\newcommand{\esempio}[1]{\ex{}{#1}}
\newcommand{\teorema}[1]{\thm{}{#1}}
\newcommand{\dimostrazione}[1]{\pf{Dimostrazione}{#1}}
\setlength{\parindent}{0pt}
\newcommand{\osservazione}[1]{\nt{#1}}
\newcommand{\inside}[1]{
  \langle #1 \rangle
}
      

%From M275 "Topology" at SJSU
\newcommand{\id}{\mathrm{id}}
\newcommand{\taking}[1]{\xrightarrow{#1}}
\newcommand{\inv}{^{-1}}

%From M170 "Introduction to Graph Theory" at SJSU
\DeclareMathOperator{\diam}{diam}
\DeclareMathOperator{\ord}{ord}
\newcommand{\defeq}{\overset{\mathrm{def}}{=}}

%From the USAMO .tex files
\newcommand{\ts}{\textsuperscript}
\newcommand{\dg}{^\circ}
\newcommand{\ii}{\item}

% % From Math 55 and Math 145 at Harvard
% \newenvironment{subproof}[1][Proof]{%
% \begin{proof}[#1] \renewcommand{\qedsymbol}{$\blacksquare$}}%
% {\end{proof}}

\newcommand{\liff}{\leftrightarrow}
\newcommand{\lthen}{\rightarrow}
\newcommand{\opname}{\operatorname}
\newcommand{\surjto}{\twoheadrightarrow}
\newcommand{\injto}{\hookrightarrow}
\newcommand{\On}{\mathrm{On}} % ordinals
\DeclareMathOperator{\img}{im} % Image
\DeclareMathOperator{\Img}{Im} % Image
\DeclareMathOperator{\coker}{coker} % Cokernel
\DeclareMathOperator{\Coker}{Coker} % Cokernel
\DeclareMathOperator{\Ker}{Ker} % Kernel
\DeclareMathOperator{\rank}{rank}
\DeclareMathOperator{\Spec}{Spec} % spectrum
\DeclareMathOperator{\Tr}{Tr} % trace
\DeclareMathOperator{\pr}{pr} % projection
\DeclareMathOperator{\ext}{ext} % extension
\DeclareMathOperator{\pred}{pred} % predecessor
\DeclareMathOperator{\dom}{dom} % domain
\DeclareMathOperator{\ran}{ran} % range
\DeclareMathOperator{\Hom}{Hom} % homomorphism
\DeclareMathOperator{\Mor}{Mor} % morphisms
\DeclareMathOperator{\End}{End} % endomorphism

\newcommand{\eps}{\epsilon}
\newcommand{\veps}{\varepsilon}
\newcommand{\ol}{\overline}
\newcommand{\ul}{\underline}
\newcommand{\wt}{\widetilde}
\newcommand{\wh}{\widehat}
\newcommand{\vocab}[1]{\textbf{\color{blue} #1}}
\providecommand{\half}{\frac{1}{2}}
\newcommand{\dang}{\measuredangle} %% Directed angle
\newcommand{\ray}[1]{\overrightarrow{#1}}
\newcommand{\seg}[1]{\overline{#1}}
\newcommand{\arc}[1]{\wideparen{#1}}
\DeclareMathOperator{\cis}{cis}
\DeclareMathOperator*{\lcm}{lcm}
\DeclareMathOperator*{\argmin}{arg min}
\DeclareMathOperator*{\argmax}{arg max}
\newcommand{\cycsum}{\sum_{\mathrm{cyc}}}
\newcommand{\symsum}{\sum_{\mathrm{sym}}}
\newcommand{\cycprod}{\prod_{\mathrm{cyc}}}
\newcommand{\symprod}{\prod_{\mathrm{sym}}}
\newcommand{\Qed}{\begin{flushright}\qed\end{flushright}}
\newcommand{\parinn}{\setlength{\parindent}{1cm}}
\newcommand{\parinf}{\setlength{\parindent}{0cm}}
% \newcommand{\norm}{\|\cdot\|}
\newcommand{\inorm}{\norm_{\infty}}
\newcommand{\opensets}{\{V_{\alpha}\}_{\alpha\in I}}
\newcommand{\oset}{V_{\alpha}}
\newcommand{\opset}[1]{V_{\alpha_{#1}}}
\newcommand{\lub}{\text{lub}}
\newcommand{\del}[2]{\frac{\partial #1}{\partial #2}}
\newcommand{\Del}[3]{\frac{\partial^{#1} #2}{\partial^{#1} #3}}
\newcommand{\deld}[2]{\dfrac{\partial #1}{\partial #2}}
\newcommand{\Deld}[3]{\dfrac{\partial^{#1} #2}{\partial^{#1} #3}}
\newcommand{\lm}{\lambda}
\newcommand{\uin}{\mathbin{\rotatebox[origin=c]{90}{$\in$}}}
\newcommand{\usubset}{\mathbin{\rotatebox[origin=c]{90}{$\subset$}}}
\newcommand{\lt}{\left}
\newcommand{\rt}{\right}
\newcommand{\bs}[1]{\boldsymbol{#1}}
\newcommand{\exs}{\exists}
\newcommand{\st}{\strut}
\newcommand{\dps}[1]{\displaystyle{#1}}

\newcommand{\sol}{\setlength{\parindent}{0cm}\textbf{\textit{Solution:}}\setlength{\parindent}{1cm} }
\newcommand{\solve}[1]{\setlength{\parindent}{0cm}\textbf{\textit{Solution: }}\setlength{\parindent}{1cm}#1 \Qed}

% Things Lie
\newcommand{\kb}{\mathfrak b}
\newcommand{\kg}{\mathfrak g}
\newcommand{\kh}{\mathfrak h}
\newcommand{\kn}{\mathfrak n}
\newcommand{\ku}{\mathfrak u}
\newcommand{\kz}{\mathfrak z}
\DeclareMathOperator{\Ext}{Ext} % Ext functor
\DeclareMathOperator{\Tor}{Tor} % Tor functor
\newcommand{\gl}{\opname{\mathfrak{gl}}} % frak gl group
\renewcommand{\sl}{\opname{\mathfrak{sl}}} % frak sl group chktex 6

% More script letters etc.
\newcommand{\SA}{\mathcal A}
\newcommand{\SB}{\mathcal B}
\newcommand{\SC}{\mathcal C}
\newcommand{\SF}{\mathcal F}
\newcommand{\SG}{\mathcal G}
\newcommand{\SH}{\mathcal H}
\newcommand{\OO}{\mathcal O}

\newcommand{\SCA}{\mathscr A}
\newcommand{\SCB}{\mathscr B}
\newcommand{\SCC}{\mathscr C}
\newcommand{\SCD}{\mathscr D}
\newcommand{\SCE}{\mathscr E}
\newcommand{\SCF}{\mathscr F}
\newcommand{\SCG}{\mathscr G}
\newcommand{\SCH}{\mathscr H}

% Mathfrak primes
\newcommand{\km}{\mathfrak m}
\newcommand{\kp}{\mathfrak p}
\newcommand{\kq}{\mathfrak q}

% number sets
\newcommand{\RR}[1][]{\ensuremath{\ifstrempty{#1}{\mathbb{R}}{\mathbb{R}^{#1}}}}
\newcommand{\NN}[1][]{\ensuremath{\ifstrempty{#1}{\mathbb{N}}{\mathbb{N}^{#1}}}}
\newcommand{\ZZ}[1][]{\ensuremath{\ifstrempty{#1}{\mathbb{Z}}{\mathbb{Z}^{#1}}}}
\newcommand{\QQ}[1][]{\ensuremath{\ifstrempty{#1}{\mathbb{Q}}{\mathbb{Q}^{#1}}}}
\newcommand{\CC}[1][]{\ensuremath{\ifstrempty{#1}{\mathbb{C}}{\mathbb{C}^{#1}}}}
\newcommand{\PP}[1][]{\ensuremath{\ifstrempty{#1}{\mathbb{P}}{\mathbb{P}^{#1}}}}
\newcommand{\HH}[1][]{\ensuremath{\ifstrempty{#1}{\mathbb{H}}{\mathbb{H}^{#1}}}}
\newcommand{\FF}[1][]{\ensuremath{\ifstrempty{#1}{\mathbb{F}}{\mathbb{F}^{#1}}}}
% expected value
\newcommand{\EE}{\ensuremath{\mathbb{E}}}
\newcommand{\charin}{\text{ char }}
\DeclareMathOperator{\sign}{sign}
\DeclareMathOperator{\Aut}{Aut}
\DeclareMathOperator{\Inn}{Inn}
\DeclareMathOperator{\Syl}{Syl}
\DeclareMathOperator{\Gal}{Gal}
\DeclareMathOperator{\GL}{GL} % General linear group
\DeclareMathOperator{\SL}{SL} % Special linear group

%---------------------------------------
% BlackBoard Math Fonts :-
%---------------------------------------

%Captital Letters
\newcommand{\bbA}{\mathbb{A}}	\newcommand{\bbB}{\mathbb{B}}
\newcommand{\bbC}{\mathbb{C}}	\newcommand{\bbD}{\mathbb{D}}
\newcommand{\bbE}{\mathbb{E}}	\newcommand{\bbF}{\mathbb{F}}
\newcommand{\bbG}{\mathbb{G}}	\newcommand{\bbH}{\mathbb{H}}
\newcommand{\bbI}{\mathbb{I}}	\newcommand{\bbJ}{\mathbb{J}}
\newcommand{\bbK}{\mathbb{K}}	\newcommand{\bbL}{\mathbb{L}}
\newcommand{\bbM}{\mathbb{M}}	\newcommand{\bbN}{\mathbb{N}}
\newcommand{\bbO}{\mathbb{O}}	\newcommand{\bbP}{\mathbb{P}}
\newcommand{\bbQ}{\mathbb{Q}}	\newcommand{\bbR}{\mathbb{R}}
\newcommand{\bbS}{\mathbb{S}}	\newcommand{\bbT}{\mathbb{T}}
\newcommand{\bbU}{\mathbb{U}}	\newcommand{\bbV}{\mathbb{V}}
\newcommand{\bbW}{\mathbb{W}}	\newcommand{\bbX}{\mathbb{X}}
\newcommand{\bbY}{\mathbb{Y}}	\newcommand{\bbZ}{\mathbb{Z}}

%---------------------------------------
% MathCal Fonts :-
%---------------------------------------

%Captital Letters
\newcommand{\mcA}{\mathcal{A}}	\newcommand{\mcB}{\mathcal{B}}
\newcommand{\mcC}{\mathcal{C}}	\newcommand{\mcD}{\mathcal{D}}
\newcommand{\mcE}{\mathcal{E}}	\newcommand{\mcF}{\mathcal{F}}
\newcommand{\mcG}{\mathcal{G}}	\newcommand{\mcH}{\mathcal{H}}
\newcommand{\mcI}{\mathcal{I}}	\newcommand{\mcJ}{\mathcal{J}}
\newcommand{\mcK}{\mathcal{K}}	\newcommand{\mcL}{\mathcal{L}}
\newcommand{\mcM}{\mathcal{M}}	\newcommand{\mcN}{\mathcal{N}}
\newcommand{\mcO}{\mathcal{O}}	\newcommand{\mcP}{\mathcal{P}}
\newcommand{\mcQ}{\mathcal{Q}}	\newcommand{\mcR}{\mathcal{R}}
\newcommand{\mcS}{\mathcal{S}}	\newcommand{\mcT}{\mathcal{T}}
\newcommand{\mcU}{\mathcal{U}}	\newcommand{\mcV}{\mathcal{V}}
\newcommand{\mcW}{\mathcal{W}}	\newcommand{\mcX}{\mathcal{X}}
\newcommand{\mcY}{\mathcal{Y}}	\newcommand{\mcZ}{\mathcal{Z}}


%---------------------------------------
% Bold Math Fonts :-
%---------------------------------------

%Captital Letters
\newcommand{\bmA}{\boldsymbol{A}}	\newcommand{\bmB}{\boldsymbol{B}}
\newcommand{\bmC}{\boldsymbol{C}}	\newcommand{\bmD}{\boldsymbol{D}}
\newcommand{\bmE}{\boldsymbol{E}}	\newcommand{\bmF}{\boldsymbol{F}}
\newcommand{\bmG}{\boldsymbol{G}}	\newcommand{\bmH}{\boldsymbol{H}}
\newcommand{\bmI}{\boldsymbol{I}}	\newcommand{\bmJ}{\boldsymbol{J}}
\newcommand{\bmK}{\boldsymbol{K}}	\newcommand{\bmL}{\boldsymbol{L}}
\newcommand{\bmM}{\boldsymbol{M}}	\newcommand{\bmN}{\boldsymbol{N}}
\newcommand{\bmO}{\boldsymbol{O}}	\newcommand{\bmP}{\boldsymbol{P}}
\newcommand{\bmQ}{\boldsymbol{Q}}	\newcommand{\bmR}{\boldsymbol{R}}
\newcommand{\bmS}{\boldsymbol{S}}	\newcommand{\bmT}{\boldsymbol{T}}
\newcommand{\bmU}{\boldsymbol{U}}	\newcommand{\bmV}{\boldsymbol{V}}
\newcommand{\bmW}{\boldsymbol{W}}	\newcommand{\bmX}{\boldsymbol{X}}
\newcommand{\bmY}{\boldsymbol{Y}}	\newcommand{\bmZ}{\boldsymbol{Z}}
%Small Letters
\newcommand{\bma}{\boldsymbol{a}}	\newcommand{\bmb}{\boldsymbol{b}}
\newcommand{\bmc}{\boldsymbol{c}}	\newcommand{\bmd}{\boldsymbol{d}}
\newcommand{\bme}{\boldsymbol{e}}	\newcommand{\bmf}{\boldsymbol{f}}
\newcommand{\bmg}{\boldsymbol{g}}	\newcommand{\bmh}{\boldsymbol{h}}
\newcommand{\bmi}{\boldsymbol{i}}	\newcommand{\bmj}{\boldsymbol{j}}
\newcommand{\bmk}{\boldsymbol{k}}	\newcommand{\bml}{\boldsymbol{l}}
\newcommand{\bmm}{\boldsymbol{m}}	\newcommand{\bmn}{\boldsymbol{n}}
\newcommand{\bmo}{\boldsymbol{o}}	\newcommand{\bmp}{\boldsymbol{p}}
\newcommand{\bmq}{\boldsymbol{q}}	\newcommand{\bmr}{\boldsymbol{r}}
\newcommand{\bms}{\boldsymbol{s}}	\newcommand{\bmt}{\boldsymbol{t}}
\newcommand{\bmu}{\boldsymbol{u}}	\newcommand{\bmv}{\boldsymbol{v}}
\newcommand{\bmw}{\boldsymbol{w}}	\newcommand{\bmx}{\boldsymbol{x}}
\newcommand{\bmy}{\boldsymbol{y}}	\newcommand{\bmz}{\boldsymbol{z}}

%---------------------------------------
% Scr Math Fonts :-
%---------------------------------------

\newcommand{\sA}{{\mathscr{A}}}   \newcommand{\sB}{{\mathscr{B}}}
\newcommand{\sC}{{\mathscr{C}}}   \newcommand{\sD}{{\mathscr{D}}}
\newcommand{\sE}{{\mathscr{E}}}   \newcommand{\sF}{{\mathscr{F}}}
\newcommand{\sG}{{\mathscr{G}}}   \newcommand{\sH}{{\mathscr{H}}}
\newcommand{\sI}{{\mathscr{I}}}   \newcommand{\sJ}{{\mathscr{J}}}
\newcommand{\sK}{{\mathscr{K}}}   \newcommand{\sL}{{\mathscr{L}}}
\newcommand{\sM}{{\mathscr{M}}}   \newcommand{\sN}{{\mathscr{N}}}
\newcommand{\sO}{{\mathscr{O}}}   \newcommand{\sP}{{\mathscr{P}}}
\newcommand{\sQ}{{\mathscr{Q}}}   \newcommand{\sR}{{\mathscr{R}}}
\newcommand{\sS}{{\mathscr{S}}}   \newcommand{\sT}{{\mathscr{T}}}
\newcommand{\sU}{{\mathscr{U}}}   \newcommand{\sV}{{\mathscr{V}}}
\newcommand{\sW}{{\mathscr{W}}}   \newcommand{\sX}{{\mathscr{X}}}
\newcommand{\sY}{{\mathscr{Y}}}   \newcommand{\sZ}{{\mathscr{Z}}}


%---------------------------------------
% Math Fraktur Font
%---------------------------------------

%Captital Letters
\newcommand{\mfA}{\mathfrak{A}}	\newcommand{\mfB}{\mathfrak{B}}
\newcommand{\mfC}{\mathfrak{C}}	\newcommand{\mfD}{\mathfrak{D}}
\newcommand{\mfE}{\mathfrak{E}}	\newcommand{\mfF}{\mathfrak{F}}
\newcommand{\mfG}{\mathfrak{G}}	\newcommand{\mfH}{\mathfrak{H}}
\newcommand{\mfI}{\mathfrak{I}}	\newcommand{\mfJ}{\mathfrak{J}}
\newcommand{\mfK}{\mathfrak{K}}	\newcommand{\mfL}{\mathfrak{L}}
\newcommand{\mfM}{\mathfrak{M}}	\newcommand{\mfN}{\mathfrak{N}}
\newcommand{\mfO}{\mathfrak{O}}	\newcommand{\mfP}{\mathfrak{P}}
\newcommand{\mfQ}{\mathfrak{Q}}	\newcommand{\mfR}{\mathfrak{R}}
\newcommand{\mfS}{\mathfrak{S}}	\newcommand{\mfT}{\mathfrak{T}}
\newcommand{\mfU}{\mathfrak{U}}	\newcommand{\mfV}{\mathfrak{V}}
\newcommand{\mfW}{\mathfrak{W}}	\newcommand{\mfX}{\mathfrak{X}}
\newcommand{\mfY}{\mathfrak{Y}}	\newcommand{\mfZ}{\mathfrak{Z}}
%Small Letters
\newcommand{\mfa}{\mathfrak{a}}	\newcommand{\mfb}{\mathfrak{b}}
\newcommand{\mfc}{\mathfrak{c}}	\newcommand{\mfd}{\mathfrak{d}}
\newcommand{\mfe}{\mathfrak{e}}	\newcommand{\mff}{\mathfrak{f}}
\newcommand{\mfg}{\mathfrak{g}}	\newcommand{\mfh}{\mathfrak{h}}
\newcommand{\mfi}{\mathfrak{i}}	\newcommand{\mfj}{\mathfrak{j}}
\newcommand{\mfk}{\mathfrak{k}}	\newcommand{\mfl}{\mathfrak{l}}
\newcommand{\mfm}{\mathfrak{m}}	\newcommand{\mfn}{\mathfrak{n}}
\newcommand{\mfo}{\mathfrak{o}}	\newcommand{\mfp}{\mathfrak{p}}
\newcommand{\mfq}{\mathfrak{q}}	\newcommand{\mfr}{\mathfrak{r}}
\newcommand{\mfs}{\mathfrak{s}}	\newcommand{\mft}{\mathfrak{t}}
\newcommand{\mfu}{\mathfrak{u}}	\newcommand{\mfv}{\mathfrak{v}}
\newcommand{\mfw}{\mathfrak{w}}	\newcommand{\mfx}{\mathfrak{x}}
\newcommand{\mfy}{\mathfrak{y}}	\newcommand{\mfz}{\mathfrak{z}}


\usepackage[T1]{fontenc}
\usepackage[utf8]{inputenc}
\usepackage{xcolor}
\usepackage{listings}
\usepackage{array}
\usepackage{longtable}

%==============================================================
% Stile per i blocchi di codice
%==============================================================
\definecolor{codebg}{HTML}{F7F7FB}
\definecolor{codekw}{HTML}{2D6FB1}
\definecolor{codecm}{HTML}{388C46}
\definecolor{codestr}{HTML}{C74440}
\definecolor{coderule}{HTML}{CCCCCC}

\lstset{
    basicstyle=\ttfamily\footnotesize,
    keywordstyle=\color{codekw}\bfseries,
    commentstyle=\color{codecm}\itshape,
    stringstyle=\color{codestr},
    numberstyle=\tiny\color{gray},
    showstringspaces=false,
    breaklines=true,
    breakatwhitespace=false,
    frame=single,
    rulecolor=\color{coderule},
    backgroundcolor=\color{codebg},
    tabsize=2,
    columns=fullflexible,
    keepspaces=true,
    captionpos=b,
    xleftmargin=1.2em,
    aboveskip=0.9em,
    belowskip=0.9em,
}

\newcommand{\code}[1]{\texttt{\small #1}}
\newcommand{\fonteslide}[1]{{\footnotesize\itshape\color{gray} Fonte: slide \emph{TemplatesInCpp}, pag.~#1.\par}}
\newcommand{\fontewiki}[1][Template (C++)]{{\footnotesize\itshape\color{gray} Fonte: Wikipedia, \emph{#1}.\par}}
\newcommand{\fontelibro}[1]{{\footnotesize\itshape\color{gray} Fonte: \emph{Thinking in C++} (Eckel), Vol.~2, #1.\par}}
\newcommand{\fonteprove}[1]{{\footnotesize\itshape\color{gray} Verificato in \code{prove/#1} (compilato ed eseguito).\par}}
\newcommand{\notaext}{{\footnotesize\itshape\color{gray} Nota: approfondimento non tratto da slide/Wikipedia, per completare il quadro.\par}}

\graphicspath{{./img/}}

%==============================================================
\begin{document}

\begin{titlepage}
    \centering
    \vspace*{2.5cm}
    {\Huge\sffamily\bfseries Template in C++\par}
    \vspace{0.5cm}
    {\LARGE\sffamily Teoria completa\par}
    \vspace{1.4cm}
    {\Large Appunti di Ingegneria del Software\par}
    \vspace{0.4cm}
    {\large Basato sulle slide del corso (\emph{Templates in C++}) e sulla voce\\
    Wikipedia \emph{Template (C++)}\par}
    \vspace{2cm}
    {\large\itshape Function template, class template, specializzazione,\\
    template variadici, template metaprogramming, concepts\par}
    \vfill
    {\large Anno 3\par}
\end{titlepage}

\tableofcontents
\newpage

%==============================================================
\chapter{Perche' servono i template}\label{ch:perche}
%==============================================================

\section{Il problema: duplicazione di codice}

Molti programmi usano strutture dati comuni: stack, code, liste. Il problema
di partenza, cosi' come lo pone il corso, e' molto concreto: un programma ha
bisogno di una \textbf{coda di clienti}; poco dopo serve anche una
\textbf{coda di messaggi}. La tentazione naturale e' prendere il codice della
coda di clienti e, con un copia-incolla-trova-sostituisci, produrre una coda
di messaggi. Quando il programma cresce ancora e serve una \textbf{coda di
ordini}, la stessa operazione va ripetuta un'altra volta.

\fonteslide{3}

\nt{
    Il problema non e' scrivere la prima coda: e' che il codice viene
    \emph{duplicato}. Se in seguito serve modificare l'implementazione della
    coda (correggere un bug, aggiungere un controllo), la modifica va
    replicata manualmente in ogni copia --- ed e' facile perderne una.
    Ri-inventare codice sorgente identico e' l'opposto di cio' che
    l'orientamento agli oggetti promette: il \textbf{riuso}.
}

\dfn{Type parameterization (genericita')}{
    L'idea alternativa alla duplicazione: implementare la coda in modo che
    possa contenere \textbf{un tipo arbitrario}, parametrizzando il codice
    rispetto al tipo degli elementi. Questa tecnica e' comunemente chiamata
    \textbf{generics}; in C++ si chiama \textbf{template}.
}

\fonteslide{4}

\section{Container class: le tre domande da porsi}

Le slide inquadrano il problema in modo piu' generale, parlando di
\emph{container class} riusabili (liste, stack, code, insiemi, alberi,
dizionari). Costruire container riusabili richiede di rispondere a tre
domande:

\begin{enumerate}
    \item E' possibile costruire abstraction di container \textbf{indipendenti
          dal tipo} dei loro elementi?
    \item Questi container possono mantenere \textbf{un solo tipo} di valore
          (container \emph{omogenei}) oppure \textbf{tipi diversi}
          contemporaneamente (container \emph{eterogenei})?
    \item E' possibile dare agli utenti del container accesso agli elementi
          \textbf{senza rimuoverli} e \textbf{senza esporre} i dettagli
          implementativi del container?
\end{enumerate}

\fonteslide{6}

\subsection{Soluzioni pre-OO}

In un linguaggio a \textbf{tipizzazione dinamica}, sia le abstraction
omogenee sia quelle eterogenee si costruiscono con relativa facilita'. Nei
linguaggi a \textbf{tipizzazione statica} pre-OO, invece:

\begin{itemize}
    \item il tipo degli elementi del container e' \emph{fisso}, quindi il
          riuso e' limitato;
    \item \textbf{variant record} o \textbf{union} aiutano fino a un certo
          punto a costruire container eterogenei, ma bisogna convivere con un
          \textbf{numero finito di alternative} e scrivere codice di test e
          cast "sporco" per estrarre i valori;
    \item nascondere i dettagli implementativi e' difficile: per iterare sul
          container serve introdurre una variabile che, con il suo tipo,
          rivela la struttura interna del container.
\end{itemize}

\fonteslide{7}

\subsection{Soluzioni OO}

Nei linguaggi OO a tipizzazione statica si possono costruire container sia
omogenei sia eterogenei, a patto che:

\begin{itemize}
    \item il linguaggio supporti la \textbf{substitutability}, con tutti i
          valori istanze di sottoclassi di una comune superclasse
          \code{general-element};
    \item gli oggetti conoscano il proprio \textbf{tipo dinamico} --- test e
          cast espliciti restano necessari per estrarre i valori, ma possono
          essere nascosti tramite sotto-classing.
\end{itemize}

Si possono introdurre \textbf{iteratori espliciti} per scorrere gli elementi
del container senza esporne la struttura implementativa interna.

\mprop{Per liste omogenee, i template SONO ``la'' soluzione}{
    Questa e' l'affermazione centrale delle slide: per costruire container
    \textbf{omogenei} in un linguaggio a tipizzazione statica, i template
    (o i generics) sono la soluzione di riferimento.
}

\fonteslide{8}

%==============================================================
\chapter{Cosa sono i template}\label{ch:cosa}
%==============================================================

\section{Definizione e scopo}

\dfn{Template}{
    I template forniscono un modo di \textbf{parametrizzare una classe o una
    funzione con un tipo}. Sono stati introdotti nei linguaggi a tipizzazione
    statica per permettere la creazione e la manipolazione di strutture dati
    generiche, come le container class riusabili.
}

\fonteslide{10}

E' il \textbf{compilatore} che genera il codice corretto per ogni
\textbf{istanziazione} del template: il corpo effettivamente eseguito non e'
generico, ma \emph{molto specifico} per quel particolare tipo.

I template sono usati estesamente nella \textbf{Standard Template Library}
(STL). Un singolo template da' origine a una \textbf{famiglia di funzioni o
classi}: tutte compilate a partire dallo stesso codice sorgente, ma che
operano su tipi diversi.

\fonteslide{11}

\nt{
    Secondo Wikipedia, dal punto di vista concettuale un template equivale,
    prima della compilazione, a \emph{copiare e incollare} il blocco di
    codice templatizzato nel punto in cui viene usato, sostituendo poi il
    parametro di template con il tipo attuale. Questa e' esattamente
    l'immagine mentale corretta per capire perche' il codice generato "non e'
    generico ma molto specifico".
}
\fontewiki

\subsection{Origini storiche e potere espressivo}

Le principali influenze storiche sui template C++ sono i \textbf{moduli
parametrizzati} del linguaggio CLU e i \textbf{generics} di Ada, che
precedono i template C++. I template C++ sono inoltre dimostrati essere
\textbf{Turing completi}: il sistema dei template, da solo, e' potente
abbastanza da poter esprimere qualunque computazione calcolabile --- proprio
questa proprieta' rende possibile il \emph{template metaprogramming}
(capitolo~\ref{ch:metaprogramming}).

\fontewiki

\section{``I template non sono fully generic''}

\wc{I template sono un mccanismo completamente generico}{
    Le slide affermano esplicitamente il contrario: \textbf{i template non
    sono fully generic}. L'obiettivo e' costruire selettivamente una classe
    assemblando solo le porzioni di codice necessarie, mantenendo \textbf{solo
    il necessario}. Per essere efficiente, questo meccanismo deve essere
    eseguito \textbf{a compile-time}, proprio per poter sopprimere la
    generazione di codice inutile --- e infatti \textbf{non c'e' type
    checking completo} alla definizione del template.
}

\fonteslide{12}

\section{Quando viene generato codice}

\mprop{Nessun codice alla definizione}{
    Nessun codice viene generato quando il compilatore incontra la
    \emph{definizione} di un template. Il compilatore esegue in quella fase
    ogni controllo possibile, ma \textbf{rimanda} il type checking completo e
    la compilazione effettiva alla fase di \textbf{istanziazione}.
}

Di conseguenza: le dichiarazioni e le definizioni dei metodi membro di una
class template devono stare tutte nello \textbf{stesso header file}.

\fonteslide{13}

\begin{Example}{La tecnicalita' declaration/definition/instantiation}{}
Il motivo concreto per cui serve questa regola emerge da un esempio a tre
file:
\begin{lstlisting}[language=C++]
// B.H  -- Declaration
template <class t>
class b {
public:
    b();
    ~b();
};

// B.CPP -- Definition
#include "B.H"
template <class t> b<t>::b()  { }
template <class t> b<t>::~b() { }

// MAIN.CPP -- Instantiation
#include "B.H"
void main() {
    b<int>   bi;
    b<float> bf;
}
\end{lstlisting}
Quando si compila \code{B.CPP}, il compilatore ha a disposizione sia le
dichiarazioni sia le definizioni: a questo punto \textbf{non} deve generare
alcuna definizione per la classe template, perche' non c'e' ancora nessuna
istanziazione. Quando invece si compila \code{MAIN.CPP} ci sono due
istanziazioni (\code{b<int>} e \code{b<float>}), ma in quel file il
compilatore ha solo le \textbf{dichiarazioni}, non le definizioni: se
dichiarazione e definizione stessero in file diversi (come qui), il linker
non troverebbe il codice generato. Ecco perche' per i template si tengono
dichiarazione e definizione nello \textbf{stesso header}.
\end{Example}

\fonteslide{14}

\section{Istanziazione: il meccanismo generale}

\dfn{Punto di istanziazione (instantiation point)}{
    Il punto nel codice dove un template viene effettivamente \textbf{usato}.
    Li' il compilatore conosce i tipi attuali corrispondenti ai parametri del
    template, e \textbf{genera il codice} per quei tipi specifici. Questo
    processo si chiama \textbf{template instantiation}; il codice generato e'
    una \emph{generated instance} del template.
}

Una classe generata da una class template si chiama \textbf{generated
class}; una funzione generata da una function template si chiama
\textbf{generated function}; un dato statico generato da un membro statico
di template si chiama \textbf{generated static data member}.

\mprop{Generazione una tantum}{
    Una generated instance viene generata \textbf{una sola volta}. Una volta
    generata, viene \textbf{riusata} a ogni successiva istanziazione con gli
    stessi parametri.
}

L'istanziazione di un template puo' avvenire in due modi:
\begin{itemize}
    \item \textbf{implicitamente}: basta usare il template con dei tipi
          concreti, es. \code{Z<int> zi;}
    \item \textbf{esplicitamente}: con la sintassi \code{template class
          Z<int>;}
\end{itemize}

\fonteslide{15}

\begin{Example}{Istanziazione implicita}{}
\begin{lstlisting}[language=C++]
template <class T> class Z {
public:
    Z()  {}
    ~Z() {}
    void f() {}
    void g() {}
};

int main() {
    Z<int>   zi;    // instantiation generates class Z<int>
    zi.f();         // and generates function Z<int>::f()
    Z<float> zf;    // instantiation generates class Z<float>
    zf.g();         // and generates function Z<float>::g()
    return 0;
}
\end{lstlisting}
Da notare: \code{Z<int>::f()} viene generata quando viene \emph{chiamata},
non semplicemente perche' esiste la classe \code{Z<int>}; \code{Z<float>::g()}
viene generata a sua volta solo quando viene chiamata --- e non
\code{Z<float>::f()} o \code{Z<int>::g()}, che nel programma non vengono
mai usate.
\end{Example}

\fonteslide{16}

\begin{Example}{Un puntatore non forza l'istanziazione}{}
\begin{lstlisting}[language=C++]
template <class T> class Z {
public:
    Z()  {} ; ~Z() {} ; void f(){} ; void g(){} ;
};
int main() {
    Z<int>*   p_zi;   // instantiation of class Z<int> NOT required
    Z<float>* p_zf;   // instantiation of class Z<float> NOT required
    return 0;
}
\end{lstlisting}
Questa volta il compilatore \textbf{non genera alcuna definizione}: non ce
n'e' bisogno. E' analogo a dichiarare un puntatore a una classe o struct
\textbf{non ancora definita} --- il compilatore deve solo sapere che
\code{Z<int>} e' un tipo, non serve conoscerne il contenuto per dichiararne
un puntatore.
\end{Example}

\fonteslide{24}

%==============================================================
\chapter{Function template}\label{ch:function}
%==============================================================

\section{Definizione e motivazione}

\dfn{Function template}{
    Permette di definire una cosiddetta \textbf{generic function}, applicabile
    ad argomenti i cui tipi possono essere specificati piu' tardi, durante la
    \emph{chiamata} alla funzione. In questo modo si informa il compilatore
    di \textbf{non} eseguire il controllo completo di consistenza dei tipi
    della definizione della funzione finche' non si incontra una chiamata: e'
    un modo di \textbf{rilassare il type checking}.
}

\fonteslide{18}

\section{Perche' non basta l'overloading}

Con l'\textbf{overloading}, il comportamento di una funzione deve essere
completamente specificato, incluso il tipo dei suoi argomenti: il
comportamento deve essere codificato \textbf{in anticipo} per ogni tipo di
parametro. Se, nel momento in cui si \emph{scrive} la funzione, il tipo degli
argomenti attuali non e' ancora noto, l'overloading semplicemente \textbf{non
si puo' usare}.

\fonteslide{19}

\section{Sintassi di dichiarazione}

Una \emph{dichiarazione} di function template e' molto simile a quella di
una funzione normale, a parte la keyword \code{template} seguita dalla lista
dei parametri di template (tra \code{<>}), che deve apparire prima della
dichiarazione della funzione:

\begin{lstlisting}[language=C++]
template <typename T1, typename T2>
    void f(T1 t, T2* p);
\end{lstlisting}

La keyword \code{typename} puo' essere sostituita, in questo contesto, dalla
keyword \code{class} --- ma questo uso e' \textbf{deprecato}; si trova spesso
in codice vecchio, pre-standard.

Qui si dichiara la function template \code{f}, con due argomenti \code{t} e
\code{p} di tipo \code{T1} e "puntatore a \code{T2}" rispettivamente.
\code{T1} e \code{T2} sono i \textbf{parametri di template}: agiscono come
placeholder per nomi di tipo attuali, che verranno specificati alla chiamata
della function template.

\fonteslide{20}

\section{Sintassi di definizione}

La \emph{definizione} di una function template segue la stessa sintassi
della dichiarazione per quanto riguarda il prototipo; il \textbf{corpo} della
definizione e' invece sintatticamente identico al corpo di una funzione
normale. All'interno del corpo, i parametri di template sono trattati
\textbf{esattamente come nomi di tipo}, dal punto di vista sintattico.

Il compilatore non puo' eseguire, \emph{in questa fase}, tutti i controlli
che eseguirebbe su nomi di tipo reali, perche' i parametri di template non
specificano ancora, in questa fase, un tipo vero e completamente definito.

\fonteslide{21}

\begin{Example}{Il classico esempio \code{r = 2 * t}}{}
\begin{lstlisting}[language=C++]
template <typename T1, typename T2> void f(T1 t, T2& r) {
    r = 2 * t;
}
\end{lstlisting}
Qui si dichiara una function template che esegue una semplice operazione
aritmetica. Non sapendo cosa siano realmente \code{T1} e \code{T2}, il
compilatore \textbf{non puo' controllare}, mentre analizza questo codice, se
l'espressione \code{r = 2 * t} sia legale. Sarebbe legale se \code{T1} fosse
\code{int} o \code{float} e \code{T2} fosse \code{float}; sarebbe invece
\textbf{illegale} se \code{T1} fosse un tipo array arbitrario (perche' un
array non si puo' moltiplicare per 2, come invece si tenta di fare con
\code{t}). E' quindi chiaro che alcuni controlli devono essere rimandati fino
a quando il template non viene \textbf{istanziato}.
\end{Example}

\fonteslide{22}

\section{Istanziazione delle function template}

Le function template non vengono effettivamente compilate finche' non sono
\textbf{esplicitamente istanziate}, cioe' usate per eseguire una chiamata.
Per chiamare una function template, la sintassi formale e', quasi sempre,
identica a quella delle chiamate a funzioni normali.

Il compilatore riesce \emph{di solito} a \textbf{dedurre automaticamente} i
tipi attuali (i "valori" dei parametri di template), usando algoritmi di
pattern matching piuttosto complessi, esaminando i tipi degli argomenti
effettivi fornito nella chiamata.

\fonteslide{23}

\begin{Example}{Deduzione corretta ed errore rilevato solo all'istanziazione}{}
\begin{lstlisting}[language=C++]
template <typename T1, typename T2> void f(T1 t, T2& r) {
    r = 2 * t;
}
class A {};   // empty class; for illustration purposes

int main() {
    int x = 1;
    int y = 2;
    f(y, x);     // OK
    A anA;
    f(anA, x);   // WRONG
}
\end{lstlisting}
\code{f(y, x)} \textbf{compila}: il compilatore deduce facilmente che
\code{T1} e' \code{int} e \code{T2} e' "reference to int", perche' sia
\code{y} sia \code{x} sono \code{int}. \code{f(anA, x)} invece
\textbf{non compila} (e il compilatore emette un errore proprio in questo
punto): il compilatore deduce senza problemi che \code{T1} e' \code{A} e
\code{T2} e' "reference to int" (perche' \code{anA} e' una \code{A} e
\code{x} e' un \code{int}), ma la funzione tenterebbe poi di moltiplicare
per 2 una copia di \code{anA}, operazione impossibile perche' la classe
\code{A} non sovraccarica \code{operator*} in modo adatto. E' esattamente
l'errore che, per costruzione, non poteva essere rilevato alla definizione
del template (sezione precedente): emerge solo ora, all'istanziazione.
\end{Example}

\fonteslide{25}

\begin{Example}{\code{max}: deduzione del tipo di ritorno}{}
\begin{lstlisting}[language=C++]
template <class T> T max(T a, T b) {
    return a > b ? a : b;
}
void main() {
    int I;
    I = max(10, 15);
    cout << "max(10, 15) = " << I << endl;
    char c;
    c = max('k', 's');
    cout << " max('k', 's') = " << c << endl;
}
\end{lstlisting}
Output: \code{max(10, 15) = 15} (\code{T} dedotto \code{int}); \code{max('k',
's') = s} (\code{T} dedotto \code{char}, confronto lessicografico sui codici
carattere).
\end{Example}

\fonteslide{26}

\section{Disambiguazione manuale}

Ci sono casi in cui il compilatore \textbf{non riesce} a dedurre il tipo
attuale del parametro di template dalla chiamata, perche' gli argomenti
attuali della funzione non portano abbastanza informazione di tipo. In questi
casi il compilatore emette un errore, segnalando la chiamata come
\textbf{ambigua} --- anche se la chiamata sarebbe eseguibile, conoscendo i
tipi da assegnare ai parametri di template.

E' quindi necessario \textbf{disambiguare la chiamata manualmente},
specificando i "valori" dei parametri di template in una lista (tra
\code{<>}) subito dopo il nome della funzione e prima della lista degli
argomenti:
\begin{lstlisting}[language=C++]
template_func_name<parameter_list>(arguments_list)
\end{lstlisting}

\fonteslide{27}

%==============================================================
\chapter{Class template}\label{ch:class}
%==============================================================

\section{Definizione}

\dfn{Class template}{
    Una classe \textbf{parametrizzata rispetto a uno o piu' tipi}. La sintassi
    per definire una class template e' quasi identica a quella delle function
    template.
}

\fonteslide{29}

\section{La differenza fondamentale: nessuna deduzione}

\mprop{Il compilatore non deduce i parametri di una class template}{
    La differenza principale fra function template e class template: per le
    \textbf{classi} il compilatore \textbf{non tenta} di dedurre i parametri
    di template. La deduzione sarebbe comunque \textbf{impossibile}: non
    c'e' nessuna informazione aggiuntiva su cui basarsi per dedurli (nel caso
    delle function template c'e' la lista degli argomenti attuali della
    chiamata, su cui la deduzione si appoggia).
}

Ogni volta che una class template viene istanziata, i "valori" (cioe' i tipi
attuali) dei parametri di template devono essere \textbf{specificati}.

\fonteslide{30}

\section{Specializzazioni come classi indipendenti}

\dfn{Specializzazione (di class template)}{
    Ogni volta che il compilatore istanzia una class template con parametri
    attuali, crea una nuova \textbf{specializzazione} --- se una
    specializzazione con quei parametri non esiste gia' (altrimenti viene
    riusata).
}

\mprop{Due specializzazioni sono classi scorrelate}{
    Due specializzazioni diverse della stessa class template sono classi
    \textbf{regolari}, e sono \textbf{completamente scorrelate} anche se
    generate dallo stesso template. Se una class template e' dichiarata come
    \code{A<typename T>}, allora \code{A<int>}, \code{A<double>} e
    \code{A<B>} (con \code{B} un'altra classe) sono classi completamente
    indipendenti fra loro.
}

Conseguenza pratica: \textbf{non} si puo' assegnare una variabile di tipo
\code{A<int>} a una variabile di tipo \code{A<double>}, a meno di fornire un
\textbf{operatore di conversione esplicito}.

\fonteslide{31}

\section{Istanziazione implicita ed esplicita}

\begin{center}
\renewcommand{\arraystretch}{1.3}
\begin{tabular}{|>{\raggedright\arraybackslash}p{6.3cm}|>{\raggedright\arraybackslash}p{6.3cm}|}
\hline
\textbf{implicita} & \textbf{esplicita} \\
\hline
\begin{lstlisting}[language=C++,frame=none,backgroundcolor=\color{white}]
Z<int> zi;
// implicit instantiation
// generates class Z<int>
\end{lstlisting}
&
\begin{lstlisting}[language=C++,frame=none,backgroundcolor=\color{white}]
template class Z<int>;
// explicit instantiation
// of class Z<int>
\end{lstlisting}
\\
\hline
\end{tabular}
\end{center}

\fonteslide{32}

\begin{Example}{Class template completa: \code{A<T>}}{}
\begin{lstlisting}[language=C++]
template <typename T> class A {
    T x, y;
public:
    A() { x = 0; y = 0; }
    T getSum() { return x + y; }
    void setXY(T a, T b) { x = a; y = b; }
};

int main() {
    A<int>    iA;
    A<double> dA;
    int y    = iA.getSum();
    double w = dA.getSum();
    iA.setXY(1,2);
    dA.setXY(0.0,0.2);
}
\end{lstlisting}
I parametri attuali di template devono essere specificati esplicitamente:
\code{A<int>} e \code{A<double>} sono tipi completamente diversi, che in
generale non si possono mescolare tra loro.
\end{Example}

\fonteslide{33}

\begin{Example}{Il codice che il compilatore genera davvero}{}
Per la classe template dell'esempio precedente, il compilatore genera --- a
tutti gli effetti --- due classi regolari indipendenti:
\begin{lstlisting}[language=C++]
class A_int {
    int x, y;
public:
    A_int() { x = 0; y = 0; }
    int getSum() { return x + y; }
    void setXY(int a, int b) { x = a; y = b; }
};
class A_double {
    double x, y;
public:
    A_double() { x = 0; y = 0; }
    double getSum() { return x + y; }
    void setXY(double a, double b) { x = a; y = b; }
};
int main() {
    A_int    iA;
    A_double dA;
    int    y = iA.getSum();
    double w = dA.getSum();
    iA.setXY(1,2);
    dA.setXY(0.0,0.2);
}
\end{lstlisting}
(I nomi \code{A\_int}/\code{A\_double} sono qui solo illustrativi: nel
binario reale i simboli generati sono \emph{mangled}, ma il punto delle
slide e' proprio questo --- il corpo eseguito \textbf{non e' generico}, e'
codice specifico duplicato per ciascun tipo istanziato.)
\end{Example}

\fonteslide{34}

%==============================================================
\chapter{Template specialization}\label{ch:specialization}
%==============================================================

\section{Cos'e' e perche' serve}

In alcuni casi e' possibile \textbf{sovrascrivere} il codice generato dal
template, fornendo una definizione speciale per un tipo specifico. Questo si
chiama \textbf{template specialization}.

\fonteslide{35}

\wc{Il comportamento generico del template va sempre bene per ogni tipo}{
    Le slide costruiscono apposta un controesempio per mostrare che non e'
    cosi': un algoritmo scritto in modo completamente generico puo'
    comportarsi in modo tecnicamente corretto ma \textbf{semanticamente
    sbagliato} per un tipo particolare.
}

\begin{Example}{Il tranello \code{max("Aladdin", "Jasmine")}}{}
\begin{lstlisting}[language=C++]
// max restituisce il massimo dei due elementi di tipo T, dove T e'
// una classe o un tipo per cui operator> e' definito
template <class T> T max(T a, T b) {
    return a > b ? a : b;
}
int main() {
    cout << "max(10, 15) = "     << max(10, 15)     << endl;
    cout << "max('k', 's') = "   << max('k', 's')   << endl;
    cout << "max(10.1, 15.2) = " << max(10.1, 15.2) << endl;
    cout << "max(\"Aladdin\", \"Jasmine\") = "
         << max("Aladdin", "Jasmine") << endl;
}
\end{lstlisting}
I primi tre risultati sono quelli attesi: \code{15}, \code{s}, \code{15.2}.
Il quarto \textbf{non} e' quello atteso (non stampa \code{"Jasmine"}, che
viene dopo alfabeticamente): la chiamata \code{max("Aladdin","Jasmine")} fa
istanziare il template con \code{T = char*}, cioe' \code{max(char*, char*)},
che confronta gli \textbf{indirizzi di memoria} dei due puntatori, non il
\textbf{contenuto} delle stringhe che puntano.
\end{Example}

\fonteslide{36-37}

\begin{Example}{La correzione tramite specializzazione}{}
\begin{lstlisting}[language=C++]
// Specialization of max for char*
template <> char* max(char* a, char* b) {
    return strcmp(a, b) > 0 ? a : b;
}
\end{lstlisting}
La specializzazione \textbf{sostituisce} la versione generica ogni volta che
\code{max} viene chiamata con due \code{char*}, usando \code{strcmp} per un
vero confronto del contenuto delle stringhe invece del confronto (privo di
significato logico) tra i loro indirizzi.

\fontewiki
Secondo Wikipedia, questo stesso esempio, riportato nella versione moderna
del linguaggio, usa \code{std::strcmp} per fornire un comportamento
\emph{ben definito} nel confronto di C-string --- il confronto diretto tra
puntatori sarebbe altrimenti privo di significato utile. In C++ moderno il
problema semplicemente \textbf{non si presenta} usando \code{std::string} al
posto di \code{char*}, perche' \code{std::string} sovraccarica
correttamente \code{operator>} confrontando il contenuto.
\end{Example}

\subsection{Le due forme di specializzazione (da Wikipedia)}

\dfn{Full specialization vs partial specialization}{
    \textbf{Full specialization}: tutti i parametri di template vengono
    specializzati (l'esempio di \code{max(char*, char*)} sopra e' una full
    specialization). \textbf{Partial specialization}: solo un
    \emph{sottoinsieme} dei parametri viene specializzato. La specializzazione
    parziale e' disponibile \textbf{solo per le class template}: le function
    template non possono essere specializzate parzialmente (si usa
    l'overloading al loro posto).
}

\fontewiki

%==============================================================
\chapter{Oltre le slide: funzionalita' moderne (da Wikipedia)}\label{ch:moderne}
%==============================================================

Il corso presenta i template nella loro forma "classica". La voce Wikipedia
\emph{Template (C++)} descrive anche le evoluzioni introdotte dagli standard
successivi (C++11, 14, 17, 20), utili per capire dove i template sono andati
dopo le basi viste finora. Tutto questo capitolo ha come unica fonte
Wikipedia.

\fontewiki

\section{Variable template (C++14)}

Da C++14 anche le \textbf{variabili} possono essere templatizzate:
\begin{lstlisting}[language=C++]
template <typename T>
constexpr T PI = T{3.141592653589793238462643383L};
\end{lstlisting}
\code{PI<float>} e \code{PI<double>} sono due costanti distinte, ciascuna
del tipo richiesto, generate a compile time.

\section{Parametri di template non-type}

Un template puo' essere parametrizzato anche su un \textbf{valore}, non solo
su un tipo:
\begin{lstlisting}[language=C++]
template <int K>
class MyClass;
\end{lstlisting}
L'esempio reale piu' importante e' \code{std::array} della libreria
standard:
\begin{lstlisting}[language=C++]
template <class T, size_t N>
struct Array;

Array<char, 6> myArray;   // array di 6 caratteri
\end{lstlisting}
\code{std::array} e' parametrizzato sia su un tipo (\code{T}) sia su un
valore intero (\code{N}, la dimensione): a differenza di un array C
classico, la dimensione fa parte del \textbf{tipo} dell'oggetto.

\section{Template variadici (C++11)}

I template variadici accettano un \textbf{numero variabile} di parametri di
template, in modo analogo alle funzioni variadiche in stile \code{printf}.

\dfn{Vantaggio principale}{
    In C++ sono supportati solo i parametri variadici "in stile C"
    (\code{...}); l'unico modo per ottenere funzioni variadiche
    \textbf{type-safe} (come in Java) e' proprio tramite i template
    variadici.
}

\begin{lstlisting}[language=C++]
template <typename... Args>
void log(const format_string<Args...>& fmt, Args&&... args) {
    std::println(logFile, fmt, args...);
}
\end{lstlisting}

\section{Template alias (C++11)}

Un template alias funziona come un \textbf{typedef parametrizzato},
utile per dare un nome comodo a un'istanziazione parziale di un template
piu' complesso:
\begin{lstlisting}[language=C++]
template <typename K, typename V>
using TreeMap = std::map<K, V>;

template <typename K, typename V>
using HashMap = std::unordered_map<K, V>;

// si puo' anche fissare solo alcuni parametri:
template <typename K>
using StringHashMap = HashMap<K, std::string>;

StringHashMap<int> myMap;   // == std::unordered_map<int, std::string>
\end{lstlisting}

\section{Function template abbreviati e concepts (C++20)}

Da C++20, usare \code{auto} (eventualmente vincolato da un \emph{concept})
come tipo di un parametro di funzione dichiara implicitamente una function
template abbreviata:
\begin{lstlisting}[language=C++]
void f1(auto x);                 // == template <typename T> void f1(T x);
void f2(Concept1 auto x);        // == template <Concept1 T> void f2(T x);
void f3(Concept2 auto... xs);    // variadico, vincolato da Concept2
\end{lstlisting}

\dfn{Concept}{
    Un \textbf{predicato booleano valutato a compile time}, usato per
    vincolare i parametri di template. Sono l'equivalente C++ dei wildcard
    di Java o delle clausole \code{where} di Rust.
}

\begin{lstlisting}[language=C++]
template <std::derived_from<Player> T>
void processListOfPlayers(const std::vector<T>& players) { /* ... */ }
\end{lstlisting}
Qui il compilatore \textbf{rifiuta} qualunque \code{T} che non derivi da
\code{Player} --- un vincolo espresso direttamente nella dichiarazione del
template, invece che scoperto tardi tramite un errore di sostituzione
oscuro (come succedeva, storicamente, con SFINAE).

\section{Template metaprogramming}\label{ch:metaprogramming}

\dfn{Template metaprogramming}{
    Usare il meccanismo dei template per eseguire \textbf{computazioni a
    compile time}, sfruttando il fatto che i template C++ sono Turing
    completi (capitolo~\ref{ch:cosa}).
}

\begin{Example}{Fattoriale a compile time, stile classico (specializzazione)}{}
\begin{lstlisting}[language=C++]
// caso induttivo
template <unsigned int N>
struct Factorial {
    static constexpr unsigned int value = N * Factorial<N - 1>::value;
};
// caso base, tramite full specialization
template <>
struct Factorial<0> {
    static constexpr unsigned int value = 1;
};
// uso: Factorial<6>::value calcola 6! a compile time
\end{lstlisting}
Questo e' lo stile "storico" di metaprogramming: si usa la
\textbf{specializzazione} (capitolo~\ref{ch:specialization}) come se fosse un
caso base di una ricorsione, esattamente come si farebbe a runtime con una
funzione ricorsiva.
\end{Example}

\begin{Example}{Lo stesso fattoriale, stile moderno (\code{constexpr} + \code{if constexpr})}{}
\begin{lstlisting}[language=C++]
template <unsigned int N>
constexpr unsigned int factorial() noexcept {
    if constexpr (N <= 1) {
        return 1;
    } else {
        return N * factorial<N - 1>();
    }
}
\end{lstlisting}
Con \code{constexpr} e \code{if constexpr} (C++17) si ottiene lo stesso
risultato con una sintassi molto piu' simile a una normale funzione
ricorsiva, senza dover simulare l'\code{if} tramite la specializzazione.
Per questo motivo oggi il template metaprogramming "puro" e' usato
soprattutto per fare operazioni \textbf{sui tipi} (non sui valori), dove
\code{constexpr} non basta.
\end{Example}

\section{Una parentesi storica: gli ``exported template''}

C++03 introdusse la keyword \code{export} per permettere di dichiarare un
template in un header senza doverne fornire la definizione in ogni unita' di
traduzione che lo usa (il vincolo visto nel capitolo~\ref{ch:cosa}). La
feature venne \textbf{rimossa in C++11}, perche' quasi nessun compilatore la
implementava davvero (solo Comeau C/C++), era costosa/complessa da
implementare, dava scarso beneficio pratico, e limitava l'evoluzione futura
del linguaggio.

Con C++20 e i \textbf{moduli}, la keyword \code{export} e' stata
reintrodotta, ma con un significato diverso (esportare simboli da un modulo):
il beneficio di velocita' di compilazione che gli exported template
avrebbero dovuto offrire e' oggi fornito dai moduli stessi, rendendo la
vecchia feature ``essenzialmente obsoleta''.

\fontewiki

%==============================================================
\chapter{Genericita' in Java}\label{ch:java}
%==============================================================

Il corso non ha slide dedicate ai generics Java: questo intero capitolo si
appoggia alla voce Wikipedia \emph{Generics in Java}.

\fontewiki[Generics in Java]

\section{Motivazione storica}

I generics sono stati introdotti in Java nel \textbf{2004}, con J2SE 5.0. Nel
1998 Gilad Bracha, Martin Odersky, David Stoutamire e Philip Wadler avevano
creato \emph{Generic Java}, un'estensione del linguaggio con supporto per
tipi generici; quel lavoro e' stato poi incorporato in Java, con l'aggiunta
del supporto ai wildcard. La motivazione primaria era rendere
\textbf{rilevabili a compile time} disallineamenti di tipo che prima
emergevano solo a runtime.

\begin{Example}{Il problema che i generics risolvono}{}
\begin{lstlisting}[language=Java]
List v = new ArrayList();
v.add("test");
Integer i = (Integer) v.get(0);   // ClassCastException a runtime
\end{lstlisting}
Il programmatore inserisce una \code{String} ma tenta un cast a
\code{Integer}: un errore logico che, senza generics, non viene rilevato
a compile time.
\begin{lstlisting}[language=Java]
List<String> v = new ArrayList<String>();
v.add("test");
Integer i = (Integer) v.get(0);   // ERRORE rilevato a compile time
\end{lstlisting}
Il parametro di tipo tra parentesi angolari dichiara che la lista contiene
\emph{esclusivamente} \code{String}: il compilatore puo' ora verificare la
correttezza dei tipi.
\end{Example}

\section{Le quattro forme di genericita' (dalla JLS)}

Secondo la Java Language Specification, la genericita' coinvolge quattro
elementi:

\begin{itemize}
    \item \textbf{Generic class}: una classe che dichiara una o piu' variabili
          di tipo, che agiscono da parametri;
    \item \textbf{Generic interface}: analogo per le interfacce;
    \item \textbf{Generic method}: un metodo che dichiara variabili di tipo
          come parametri di tipo formali, con sintassi analoga a quella di
          classi/interfacce generiche;
    \item \textbf{Generic constructor}: un costruttore che dichiara
          indipendentemente le proprie variabili di tipo, permettendo un
          comportamento generico anche se la classe che lo contiene non e'
          generica.
\end{itemize}

\begin{Example}{Una classe generica \code{Entry<K,V>}}{}
\begin{lstlisting}[language=Java]
public class Entry<K, V> {
    private K key;
    private V value;
    public Entry(K key, V value) { this.key = key; this.value = value; }
    public K getKey()   { return key; }
    public V getValue() { return value; }
    public String toString() { return String.format("(%s, %s)", key, value); }
}

Entry<String, String>  grade = new Entry<>("Mike", "A");
Entry<String, Integer> mark  = new Entry<>("Mike", 100);
Entry<Integer, Boolean> prime = new Entry<>(13, true);
\end{lstlisting}
E un metodo generico:
\begin{lstlisting}[language=Java]
public static <T> Entry<T, T> twice(T value) {
    return new Entry<T, T>(value, value);
}

Entry<String, String> pair = Entry.twice("Hello");          // dedotto
Entry<String, String> pair2 = Entry.<String>twice("Hello");  // esplicito
\end{lstlisting}
\end{Example}

\subsection{Diamond operator e varargs generici}

Da Java SE 7, l'\textbf{operatore diamond} (\code{<>}) permette di omettere i
parametri di tipo ridondanti quando sono deducibili dal contesto (vedi
\code{new Entry<>(...)} sopra).

\wc{Il diamond operator Java funziona come \code{Box<>} in C++}{
    No: in C++ una sintassi simile a \code{Box<>} si riferisce a
    \textbf{default template argument} (es. \code{Box<int>} con un default),
    non a inferenza di tipo. Il diamond operator Java e' invece puro
    \textbf{type inference} al punto di chiamata. Anche C\# e Kotlin non
    hanno affatto il diamond operator: usano sintassi alternativa per
    evitare la ripetizione ridondante dei tipi.
}

I metodi generici supportano anche i \textbf{varargs}:
\begin{lstlisting}[language=Java]
public <T> T[] toArray(T... elements) { return elements; }
Integer[] array = toArray(1, 2, 3, 4, 5, 6);
\end{lstlisting}

\subsection{Restrizione sui tipi primitivi}

I tipi primitivi sono \textbf{vietati} come parametri di tipo: serve la
classe wrapper. \code{Entry<int, int>} non compila; serve
\code{Entry<Integer, Integer>}. Gli \emph{array} di primitivi, essendo
oggetti essi stessi, sono invece ammessi: \code{List<int[]>} e' legale.

\section{Wildcard}

I wildcard (\code{?}) sono argomenti di tipo per tipi parametrizzati,
opzionalmente vincolati.

\begin{Example}{Le tre forme di wildcard}{}
\begin{lstlisting}[language=Java]
// unbounded: puo' leggere solo Object, scrivere solo null
Collection<?> c = new ArrayList<String>();
c.add(new Object());   // ERRORE a compile time
c.add(null);           // permesso

// upper-bounded (extends): "produttore", si legge
List<? extends Number> nums = new ArrayList<Float>();
nums.add(3.14);         // ERRORE
nums.add(null);          // permesso

// lower-bounded (super): "consumatore", si scrive
List<? super Number> nums2 = new ArrayList<Object>();
nums2.add(42);           // OK: 42 e' un Number
\end{lstlisting}
\end{Example}

\dfn{PECS (Producer Extends, Consumer Super)}{
    Mnemonico di Joshua Bloch (\emph{Effective Java}): usa \code{? extends}
    quando il parametrizzato \textbf{produce} valori (li leggi); usa
    \code{? super} quando \textbf{consuma} valori (li scrivi). Corrisponde
    esattamente alle nozioni di covarianza e controvarianza.
}

\section{Covarianza degli array vs invarianza dei generics}

\begin{Example}{Perche' i generics sono invarianti (a differenza degli array)}{}
\begin{lstlisting}[language=Java]
Object[] objs = new Long[0];
objs[0] = "string";     // compila; ArrayStoreException A RUNTIME

Collection<Object> objs2 = new ArrayList<Long>();   // ERRORE a compile time
\end{lstlisting}
Gli array in Java sono \textbf{covarianti} (\code{Long[]} e' assegnabile a
\code{Object[]}), ma questo costringe a un controllo a runtime che puo'
fallire con \code{ArrayStoreException}. I generics sono invece
deliberatamente \textbf{invarianti}: \code{ArrayList<Long>} non e' un
\code{Collection<Object>}, e l'errore viene spostato interamente a compile
time, prima ancora di eseguire il programma.
\end{Example}

\section{Type erasure: il meccanismo}

\dfn{Type erasure}{
    L'informazione di tipo generico viene \textbf{rimossa durante la
    compilazione}. \code{List<Integer>} diventa il semplice \code{List} (non
    generico) nel bytecode compilato; la correttezza e' garantita dai
    controlli fatti a compile time, non da alcuna verifica residua a
    runtime.
}

\begin{lstlisting}[language=Java]
List<Integer> li = new ArrayList<>();
List<Float>   lf = new ArrayList<>();
if (li.getClass() == lf.getClass()) {
    System.out.println("Equal");   // stampa "Equal": a runtime sono la stessa classe
}
\end{lstlisting}

\subsection{Conseguenze dell'erasure}

\begin{itemize}
    \item \textbf{Nessuna classe eccezione generica}: una classe generica non
          puo' estendere \code{Throwable}. Per erasure, un
          \code{catch (GenericException<Integer> e)} e un
          \code{catch (GenericException<String> e)} sarebbero
          indistinguibili --- il secondo sarebbe codice \textbf{irraggiungibile},
          quindi il linguaggio vieta la cosa a monte;
    \item \textbf{Nessuna creazione diretta di istanze}: \code{new T()} non
          compila, perche' a runtime \code{T} non esiste piu';
    \item \textbf{Restrizioni sui membri static}: esiste \textbf{una sola}
          copia a runtime di una classe generica; tutte le istanze
          condividono le variabili static indipendentemente dal parametro di
          tipo usato, quindi un parametro di tipo non puo' apparire in una
          dichiarazione static;
    \item nei \code{throws}, un parametro di tipo e' ammesso solo se
          \textbf{vincolato} da \code{Throwable}:
\begin{lstlisting}[language=Java]
public <T extends Throwable> void throwMeConditional(boolean c, T exception) throws T {
    if (c) throw exception;
}
\end{lstlisting}
\end{itemize}

\dfn{Reificazione}{
    Gli \textbf{array sono reificati}: portano l'informazione di tipo a
    runtime (ogni tipo \code{T} ha un corrispondente \code{T[]}). I
    \textbf{generics non sono reificati}: un oggetto di tipo generico perde,
    a runtime, i dettagli di rappresentazione che aveva a compile time. E'
    esattamente questa distinzione che spiega perche' non si possono creare
    array generici direttamente, e perche' \code{instanceof} su un tipo
    parametrizzato concreto sia vietato.
}

\subsection{Perche' l'erasure: retrocompatibilita'}

La scelta della type erasure ha mantenuto la \textbf{compatibilita'} con i
programmi Java pre-generics (pre J2SE 5.0), permettendo l'introduzione dei
generics senza richiedere modifiche alla JVM.

\section{Confronto esplicito con i template C++ (da Wikipedia)}

\begin{center}
\renewcommand{\arraystretch}{1.4}
\footnotesize
\begin{tabular}{|>{\raggedright\arraybackslash}p{4.2cm}|>{\raggedright\arraybackslash}p{4.2cm}|>{\raggedright\arraybackslash}p{4.2cm}|}
\hline
 & \textbf{C++ template} & \textbf{Java generics} \\
\hline
Codice compilato & una versione \textbf{per ogni} istanziazione di tipo & \textbf{una sola} versione compilata, indipendentemente dai parametri di tipo usati \\
\hline
Specializzazione & si', supportata (cap.~\ref{ch:specialization}) & \textbf{non} supportata \\
\hline
Informazione di tipo & disponibile a compile time \emph{e} usata per generare codice specifico & disponibile solo a compile time, poi \textbf{cancellata} \\
\hline
Sintassi ``diamond'' & \code{Box<>} = default template argument & \code{Box<>} = inferenza di tipo (Java 7+) \\
\hline
\end{tabular}
\end{center}

\fontewiki[Generics in Java]

\section{Sviluppi futuri: Project Valhalla}

Il progetto \textbf{Valhalla} studia possibili miglioramenti futuri ai
generics Java, tra cui: \textbf{generic specialization} (permettere
\code{List<int>} invece di richiedere \code{List<Integer>}) e
\textbf{reified generics} (rendere i tipi attuali disponibili a runtime).
Questi miglioramenti dovrebbero convivere con la retrocompatibilita'.

\fontewiki[Generics in Java]

%==============================================================
\chapter{Genericita' in Python}\label{ch:python}
%==============================================================

\section{Il punto di vista di Wikipedia: genericita' ``gratuita''}

La voce Wikipedia \emph{Generic programming} osserva che nei linguaggi a
\textbf{tipizzazione dinamica} il passaggio di valori alle funzioni e
l'assegnazione di valori sono \textbf{type-indifferent}: la genericita' esiste
gia' per costruzione. Questo di solito \textbf{non viene nemmeno chiamato
``genericita'''}, perche' e' una conseguenza diretta del sistema di tipi
dinamico, non un meccanismo dedicato come i template o i generics.

\fontewiki[Generic programming]

Questo e' esattamente il punto della sezione ``Container Classes Case
Study'' delle slide (capitolo~\ref{ch:perche}), dove si diceva che in un
linguaggio a tipizzazione dinamica le abstraction omogenee ed eterogenee "si
costruiscono con relativa facilita'": le slide parlano di C++ e Java, ma la
stessa osservazione vale per Python.

\section{Duck typing}

\dfn{Duck typing (duck test)}{
    "Se cammina come un'anatra e starnazza come un'anatra, allora deve essere
    un'anatra." Un oggetto e' considerato di un certo tipo se possiede
    \textbf{tutti i metodi e le proprieta'} richiesti da quel tipo,
    indipendentemente da qualunque dichiarazione formale di quel tipo.
}

\fontewiki[Duck Typing]

La differenza rispetto alla \textbf{tipizzazione nominativa} (nominal
typing): quest'ultima verifica il tipo tramite \textbf{dichiarazione
esplicita} (es. \code{class MyBag(Sized)} in Python, o
\code{implements}/\code{extends} in Java); il duck typing valuta il tipo in
base ai \textbf{comportamenti effettivamente disponibili}. Python e'
l'esempio principale di linguaggio dove il duck typing e' nativamente
integrato (insieme, in misura diversa, a Ruby e JavaScript).

\fontewiki[Duck Typing]

\section{Il collegamento teorico con i template: lo stesso ``duck test'', due momenti diversi}

Questo e' il punto di sintesi piu' importante di tutto il documento, e viene
letteralmente dalla voce Wikipedia sul duck typing:

\mprop{Template e duck typing applicano lo stesso test, in momenti diversi}{
    Le funzioni template (C++) applicano il \emph{duck test} in un contesto
    di tipizzazione \textbf{statica}, mantenendo vantaggi e svantaggi di
    entrambi gli approcci. Il duck typing e' \textbf{piu' flessibile}, perche'
    richiede solo l'implementazione dei metodi \emph{effettivamente
    invocati a runtime}; i template richiedono invece che siano valide
    \emph{tutte} le implementazioni non provabilmente irraggiungibili a
    compile time.
}

\fontewiki[Duck Typing]

\nt{
    Questa e' la stessa idea, vista da un'altra angolazione, del capitolo
    \ref{ch:cosa}: "non c'e' type checking completo" alla definizione di un
    template, perche' il controllo e' rimandato all'istanziazione --- che e'
    l'equivalente, a compile time, del momento in cui in Python
    l'\code{AttributeError} emergerebbe a runtime se un metodo mancasse
    davvero. Template e duck typing risolvono lo stesso problema
    (``non voglio impegnarmi su un'interfaccia esplicita'') a due estremi
    opposti dell'asse compile-time/runtime.
}

\section{Structural typing}

Wikipedia distingue ulteriormente: il \textbf{structural typing} e' un
sistema di tipizzazione \emph{statico} che determina la compatibilita' in
base alla \textbf{struttura} del tipo; il duck typing e' invece
\emph{dinamico} e valuta solo la parte di struttura effettivamente usata a
runtime. Linguaggi come TypeScript, Elm e \textbf{Python} supportano il
structural typing a vari livelli.

\fontewiki[Duck Typing]

\notaext
In Python il modulo \code{typing} realizza concretamente questa idea con
\code{Protocol}: una classe puo' "soddisfare" un protocollo avendo i metodi
giusti, senza dichiarare mai un'ereditarieta' esplicita --- e' duck typing
formalizzato staticamente, verificabile da uno strumento esterno come
\code{mypy}:
\begin{lstlisting}[language=Python]
from typing import Protocol, TypeVar, Generic

class Sized(Protocol):
    def __len__(self) -> int: ...

def describe(x: Sized) -> str:
    return f"size {len(x)}"

class MyBag:
    def __len__(self): return 42

describe(MyBag())   # OK per mypy: MyBag ha __len__, non serve ereditare da Sized

T = TypeVar("T")
class Box(Generic[T]):
    def __init__(self, item: T) -> None:
        self.item = item
\end{lstlisting}

\wc{I type hint di Python (\code{Generic[T]}, \code{TypeVar}) funzionano come i generics Java}{
    No: a runtime, di default, Python \textbf{non controlla affatto} i type
    hint. \code{Box[int]("ciao")} costruisce felicemente un \code{Box} con
    dentro una stringa, senza alcun errore: i type hint finiscono solo nel
    dizionario \code{\_\_annotations\_\_} della funzione/classe, pura
    metadata. L'unico modo di intercettare una violazione e' eseguire uno
    strumento di analisi statica \textbf{offline} (\code{mypy}) prima di
    lanciare il programma --- esattamente l'opposto del controllo
    obbligatorio e automatico che accompagna i generics Java o i template
    C++.
}

%==============================================================
\chapter{Genericita' in C}\label{ch:c}
%==============================================================

\section{``Function overloading, non true generic programming''}

La voce Wikipedia \emph{Generic programming} e' molto netta sul C: lo
standard C11 offre le cosiddette \textbf{type-generic expression} tramite la
keyword \code{\_Generic}, che fornisce una forma di \textbf{overloading di
funzione}, ma \textbf{non} vera programmazione generica.

\fontewiki[Generic programming]

Questo conferma, con l'autorita' di una fonte esterna al corso, esattamente
il limite discusso a proposito di \code{\_Generic}: la selezione avviene a
compile time fra un insieme \textbf{fisso e scritto a mano} di funzioni
gia' esistenti (una per tipo elencato nella macro), non generazione
automatica di codice per un tipo arbitrario futuro come farebbe un template.

\section{L'esempio GNU C: macro + \code{typeof} + statement expression}

Wikipedia mostra un esempio interessante di macro ``generica'' che risolve
un problema classico:

\begin{lstlisting}[language=C]
#define max(a,b) \
        ({ typeof (a) _a = (a); \
           typeof (b) _b = (b); \
           _a > _b ? _a : _b; })
\end{lstlisting}

\fontewiki[Generic programming]

\notaext
Questa macro e' istruttiva soprattutto per \textbf{cosa evita}: la
classica macro \code{\#define MAX(a,b) ((a)>(b)?(a):(b))} valuta i suoi
argomenti \textbf{piu' volte} per sostituzione testuale (se \code{a} e'
un'espressione con effetti collaterali come \code{i++}, viene eseguita due
volte). Qui invece \code{\_a} e \code{\_b} sono variabili temporanee, create
una sola volta dentro uno \emph{statement expression} \code{(\{ ... \})}
--- un'estensione GNU, non C standard --- e \code{a}/\code{b} vengono quindi
valutati esattamente \textbf{una volta} ciascuno. \code{typeof} (altra
estensione GNU) deduce il tipo di \code{\_a}/\code{\_b} dall'espressione
stessa, dando un minimo di "genericita' sui tipi" pur restando pura
espansione testuale del preprocessore --- senza alcun controllo di tipo
reale al momento della scrittura della macro.

\section{Il quadro completo delle tecniche C (slide + Wikipedia)}

Il capitolo sulle slide C++ (sezione~\ref{ch:perche}, "Pre OO Solutions for
Container Classes") aveva gia' indicato, per i linguaggi a tipizzazione
statica pre-OO come il C, le \textbf{variant record}/\textbf{union} come
tecnica per container eterogenei, con il costo di un numero finito di
alternative e di codice di test/cast "sporco". A quella tecnica si
affiancano, storicamente:

\begin{itemize}
    \item \textbf{macro del preprocessore} (sostituzione testuale, zero
          controllo di tipo, rischio di valutazione multipla come visto
          sopra);
    \item \textbf{\code{void*}} (puntatore generico, richiede cast manuali
          a ogni estrazione, zero type-safety --- e' la tecnica usata da
          funzioni della libreria standard come \code{qsort});
    \item \textbf{\code{\_Generic}} (C11): dispatch a compile time su un
          insieme chiuso di tipi/funzioni scritte a mano.
\end{itemize}

\wc{\code{\_Generic} e' l'equivalente C dei template}{
    Secondo la stessa Wikipedia, no: \code{\_Generic} da' overloading su un
    insieme \textbf{chiuso e predefinito} di funzioni, mentre un template
    (o un generic Java) si applica a \textbf{qualunque} tipo che soddisfi i
    requisiti richiesti, \textbf{compresi quelli che non esistevano ancora}
    quando il template e' stato scritto. Aggiungere un nuovo tipo a
    \code{\_Generic} richiede di modificare a mano la macro; aggiungere un
    nuovo tipo a un template richiede solo che quel tipo supporti le
    operazioni usate nel corpo del template.
}

%==============================================================
\chapter{Templates in Depth (dal libro di Eckel)}\label{ch:libro}
%==============================================================

Il docente ha indicato come approfondimento il libro \emph{Thinking in C++}
di Bruce Eckel, liberamente accessibile. Il Vol.~2 dedica ai template
l'intero \textbf{capitolo 5, ``Templates in Depth''}, che presenta i
template come un vero e proprio \emph{sotto-linguaggio} dentro il C++:
``i template fanno molto di piu' che creare container generici''.

\fontelibro{cap.~5}

\section{Parametri di template: le quattro forme}

\dfn{Non-type parameter}{
    Un template puo' essere parametrizzato non solo su un \emph{tipo} ma
    anche su un \textbf{valore costante noto a compile time}:
    \code{template<class T, int sz>}. Il valore diventa \textbf{parte del
    tipo}.
}

\begin{Example}{Non-type parameter, verificato}{}
\begin{lstlisting}[language=C++]
template <class T, int sz>
class Array {
    T data[sz];                  // la dimensione fa parte del TIPO
public:
    T& operator[](int i) { return data[i]; }
    int size() const { return sz; }
};

Array<int, 5>   a5;
Array<int, 100> a100;
\end{lstlisting}
Output verificato: \code{sizeof(Array<int,5>) = 20},
\code{sizeof(Array<int,100>) = 400} --- sono \textbf{tipi diversi}.
Un parametro non-type deve essere una \textbf{costante di compilazione}:
\code{int n; cin >> n; Array<int,n> a;} e' impossibile.
E' esattamente il meccanismo di \code{std::array<T,N>}.
\end{Example}

\fontelibro{cap.~5, ``Non-type parameters''}
\fonteprove{T04\_libro\_nontype\_typename.cpp}

\dfn{Default template argument}{
    Come le funzioni, i template ammettono valori di default:
    \code{template<class T = int, int sz = 10>}, cosi' che
    \code{ArrayConDefault<>} sia gia' utilizzabile.
}

\dfn{Template template parameter}{
    Un parametro che e' esso stesso un \emph{template} (non un tipo
    concreto), permettendo un'astrazione di livello superiore.
}

\mprop{Sono proprio queste le feature che mancano ai generics}{
    Default template argument, non-type argument e template template
    argument sono tre delle cinque funzionalita' che Wikipedia elenca
    come \textbf{assenti} nei generics di Java e .NET
    (capitolo~\ref{ch:confronto}).
}

\section{La keyword \code{typename} e i nomi dipendenti}

\wc{\code{typename} e \code{class} sono sempre intercambiabili}{
    Sono equivalenti \emph{solo} nella lista dei parametri di template.
    \code{typename} ha un \textbf{secondo uso, non opzionale}:
    disambiguare i \textbf{nomi dipendenti} dentro il corpo di un template.
}

\fontelibro{cap.~5, ``The typename keyword''}

\begin{Example}{Perche' serve \code{typename}}{}
\begin{lstlisting}[language=C++]
template <typename Container>
void stampaPrimo(const Container& c) {
    typename Container::const_iterator it = c.begin();
    //  ^^^^^^^^ senza questo, errore oscuro
    std::cout << *it;
}
\end{lstlisting}
Senza \code{typename}, il compilatore assume di default che
\code{Container::const\_iterator} \textbf{non sia un tipo} (potrebbe essere
un dato statico), e interpreta la riga come un'\emph{espressione}
(una moltiplicazione!). L'errore che ne risulta e' notoriamente
incomprensibile.
\end{Example}

\fonteprove{T04\_libro\_nontype\_typename.cpp}

\section{Overloading e ordinamento parziale}

Piu' function template possono coesistere: il compilatore sceglie il
\textbf{match piu' specializzato}, secondo le regole di \emph{partial
ordering}. Il libro segnala anche una complicazione pratica: prendere
l'\textbf{indirizzo} di una function template richiede un'istanziazione
esplicita, perche' il compilatore non ha argomenti da cui dedurre i tipi.

\fontelibro{cap.~5, ``Function templates''}

\section{Gli idiomi: traits, policies, CRTP}

Il capitolo dedica una sezione ai \emph{template programming idioms}, cioe'
agli usi dei template che vanno oltre il container generico.

\dfn{Traits}{
    Informazioni sul \emph{tipo} codificate dentro classi template, che
    permettono \textbf{decisioni a compile time}: proprieta' numeriche,
    categorie di iteratori, codifiche di caratteri. Consentono ad un
    algoritmo di \emph{adattarsi} al proprio parametro di template.
}

\dfn{Policies}{
    Parametri di template che \textbf{controllano il comportamento} di un
    algoritmo, permettendo un design flessibile \textbf{senza overhead a
    runtime} (l'esempio classico sono gli \emph{allocator} dei container
    della STL).
}

\dfn{CRTP --- Curiously Recurring Template Pattern}{
    Una classe \textbf{eredita da un template parametrizzato su se stessa}:
    \code{struct Der : Base<Der>}. Permette \textbf{polimorfismo statico
    senza funzioni virtuali} --- si scambia il dispatch dinamico per
    un'astrazione a costo zero.
}

\fontelibro{cap.~5, ``Template programming idioms''}

\begin{Example}{CRTP: polimorfismo senza vtable, verificato}{}
\begin{lstlisting}[language=C++]
template <typename Derived>
struct Base {
    void interface() {
        static_cast<Derived*>(this)->implementation();  // dispatch STATICO
    }
};
struct Der1 : Base<Der1> {
    void implementation() { std::cout << "Der1\n"; }
};
\end{lstlisting}
Output verificato: \code{sizeof(Der1) = 1} byte --- \textbf{nessun vptr}.
Con una funzione virtuale sarebbero stati 8 byte (il vptr).
Tutto il dispatch e' risolto a compile time e inlineabile.
\end{Example}

\fonteprove{T05\_libro\_crtp\_traits\_metaprog.cpp}

\nt{
    \textbf{Collegamento con l'altro fascicolo:} il CRTP e' l'esatto
    equivalente C++ dell'idioma Java \code{class Enum<E extends Enum<E>>}
    (\emph{F-bounded polymorphism}). Stessa idea --- una base generica che
    ``sa'' quale sia il tipo concreto che la deriva --- realizzata con due
    strategie di genericita' opposte: monomorfizzazione (nessun residuo a
    runtime) contro erasure.
}

\section{Metaprogramming ed expression templates}

Oltre al fattoriale a compile time (capitolo~\ref{ch:moderne}), il libro
introduce le \textbf{expression templates}: una tecnica sofisticata che
\emph{cattura le operazioni} dentro strutture template prima di valutarle,
permettendo ottimizzazioni --- per esempio, nell'aritmetica su vettori,
\textbf{evitare la creazione di temporanei} intermedi. E' la base di
librerie di calcolo numerico ad alte prestazioni come Eigen e Blitz++.

\fontelibro{cap.~5, ``Template metaprogramming''}

\section{I due modelli di compilazione}

\dfn{Inclusion model}{
    Le definizioni dei template stanno negli \textbf{header}. Costringe a
    ricompilare tutto quando un template cambia, ma rende immediata la
    compilazione separata. E' il modello universalmente usato, ed e' il
    motivo della regola vista nel capitolo~\ref{ch:cosa}.
}

\dfn{Separation model (\code{export})}{
    La keyword \code{export} avrebbe dovuto permettere di mettere le
    implementazioni dei template in file sorgente separati, risolvendo i
    problemi di efficienza di compilazione nei progetti grandi. Il libro la
    descrive come ``raramente implementata'' --- ed e' infatti stata
    \textbf{rimossa in C++11} (capitolo~\ref{ch:moderne}).
}

\fontelibro{cap.~5, ``Compilation models''}

\nt{
    Il libro avverte inoltre che i \textbf{nomi dentro i template} seguono
    regole di lookup complesse: i nomi dipendenti sono risolti in modo
    diverso, e le dichiarazioni \code{friend} dentro un template richiedono
    attenzione per mantenere le relazioni di accesso volute.
}

%==============================================================
\chapter{Sintesi trasversale}\label{ch:confronto}
%==============================================================

Wikipedia sottolinea un punto su cui le slide non si spingono, ma che e'
esattamente il tipo di domanda che il docente pone all'esame: template C++,
generics Java e generics .NET sono spesso considerati simili, ma
\textbf{i generics imitano solo il comportamento di base dei template C++}.

\fontewiki

\dfn{Cosa manca ai generics rispetto ai template C++}{
    Funzionalita' avanzate dei template C++ che i generics (Java, .NET) non
    offrono:
    \begin{itemize}
        \item \textbf{explicit specialization} (capitolo~\ref{ch:specialization});
        \item \textbf{partial specialization} (solo per le classi);
        \item \textbf{default template argument};
        \item \textbf{non-type template argument} (es. \code{Array<char, 6>});
        \item \textbf{template template argument}.
    \end{itemize}
    Queste funzionalita' avanzate sono cio' che rende possibili librerie
    sofisticate di template metaprogramming come Boost, STLSoft e la STL
    stessa.
}

\nt{
    Storicamente, il concetto di template non era presente in linguaggi come
    Java e C\# 1.0. Java ha poi adottato i \textbf{generics} (che imitano il
    comportamento di base dei template, ma tecnicamente sono un meccanismo
    diverso --- \emph{erasure} invece di generazione di codice); anche C\# ha
    aggiunto i generics, in .NET 2.0. I generics di Ada, al contrario,
    \textbf{precedono} i template C++.
}

\fontewiki

\section{Tabella comparativa a quattro vie}

\begin{center}
\renewcommand{\arraystretch}{1.4}
\setlength{\tabcolsep}{4pt}
\footnotesize
\begin{tabular}{|>{\raggedright\arraybackslash}p{2.0cm}|>{\raggedright\arraybackslash}p{3.0cm}|>{\raggedright\arraybackslash}p{3.0cm}|>{\raggedright\arraybackslash}p{3.0cm}|>{\raggedright\arraybackslash}p{2.6cm}|}
\hline
 & \textbf{C} & \textbf{C++} & \textbf{Java} & \textbf{Python} \\
\hline
Meccanismo & nessuno strutturato: macro, \code{void*}, \code{\_Generic} (insieme chiuso) & template: generazione di codice a compile time & generics: type erasure a compile time & nessuno: tipizzazione dinamica, duck typing \\
\hline
Quando si applica il duck test & mai in modo generale & a compile time, su chi verrebbe istanziato & a compile time, entro i bound dichiarati & a runtime, solo sui metodi davvero invocati \\
\hline
Codice a runtime & identico a quanto scritto a mano & una copia per tipo istanziato (monomorfizzazione) & una sola copia condivisa (erasure) & un solo oggetto codice, nessuna specializzazione \\
\hline
Type-safety & nessuna (cast manuali) & completa, ma solo dopo l'istanziazione & completa a compile time, persa a runtime & nessuna di default; opzionale e offline con mypy \\
\hline
\end{tabular}
\end{center}

\section{Il filo che lega tutto}

Il punto piu' importante da portare all'esame, che nessuna delle fonti
isolate dice esplicitamente ma che emerge mettendole a confronto: template
C++, generics Java e duck typing Python sono \textbf{tre risposte allo
stesso identico problema} (evitare di duplicare codice per tipi diversi,
capitolo~\ref{ch:perche}), che differiscono solo per \textbf{quando}
verificano che il tipo "vada bene" (compile time con generazione di codice,
compile time con erasure, runtime senza generazione), non per \textbf{cosa}
verificano (che il tipo supporti le operazioni richieste --- il "duck test",
capitolo~\ref{ch:python}). C, privo di entrambi i meccanismi, e' li' a
ricordare cosa succede quando quel controllo, a qualunque livello, semplicemente
\textbf{non c'e'}.

%==============================================================
\chapter{L'esercizio proposto dal corso}\label{ch:esercizio}
%==============================================================

Le slide concludono con un esercizio proposto, che chiede di completare
l'implementazione di una \code{List<T>} e di uno \code{Stack<T>} generici.

\begin{Example}{Scheletro di \code{List<T>} da completare}{}
\begin{lstlisting}[language=C++]
template <class T> class List {
public:
    void add(T);
    T firstElement;
    List<T> * nextElements;
};

List<int>    A;
List<double> B;

template <class T> int length(List<T> & X) {
    if (X==0) return 0;
    return 1 + length(X.nextElements);
}
\end{lstlisting}
\end{Example}

\fonteslide{39}

\begin{Example}{Scheletro di \code{Stack<T>} da completare}{}
\begin{lstlisting}[language=C++]
template <class T> class Stack {
public:
    void Push(T);
    T Pop();
};
template <class T> void Stack<T>::Push(T arg) {
    //...
};
template <class T> T Stack<T>::Pop() {
    //...
};

Stack<int>    s1;
Stack<double> s2;
s1.Push(1);
s1.Push(1.0);
\end{lstlisting}
\end{Example}

\fonteslide{40}

\nt{
    Da notare: la sintassi \code{template <class T> void Stack<T>::Push(T arg)
    \{ ... \}} per definire fuori dalla classe un metodo di una class
    template va sempre ripetuta con la keyword \code{template
    <class T>} anche nella definizione --- esattamente come visto nella
    tecnicalita' declaration/definition/instantiation del
    capitolo~\ref{ch:cosa}: la \emph{definizione} di \code{Push} e \code{Pop}
    e' essa stessa un template, non ancora una funzione concreta.
}

%==============================================================
\chapter{Riepilogo}\label{ch:riepilogo}
%==============================================================

\begin{center}
\renewcommand{\arraystretch}{1.4}
\footnotesize
\begin{tabular}{|>{\raggedright\arraybackslash}p{3.2cm}|>{\raggedright\arraybackslash}p{9.5cm}|}
\hline
\textbf{Concetto} & \textbf{Punto chiave} \\
\hline
Motivazione & evitare la duplicazione di codice per container che differiscono solo per il tipo degli elementi (cap.~\ref{ch:perche}) \\
\hline
Non fully generic & nessun type checking completo alla definizione; tutto e' rimandato all'istanziazione, per poter generare solo il codice necessario (cap.~\ref{ch:cosa}) \\
\hline
Function template & il compilatore \emph{deduce} i parametri dagli argomenti della chiamata; alcuni errori (es. \code{operator*} mancante) emergono solo all'istanziazione (cap.~\ref{ch:function}) \\
\hline
Class template & nessuna deduzione: i parametri vanno sempre specificati; ogni istanziazione e' una classe indipendente e scorrelata dalle altre (cap.~\ref{ch:class}) \\
\hline
Specialization & si puo' sovrascrivere il codice generato per un tipo specifico (es. \code{char*} con \code{strcmp}); parziale solo per le classi (cap.~\ref{ch:specialization}) \\
\hline
Oltre le slide & variable template, non-type/template-template parameter, variadic template, alias, concepts, metaprogramming (cap.~\ref{ch:moderne}) \\
\hline
Java & generics con \emph{type erasure}: una sola copia a runtime, wildcard e PECS, invarianza contro la covarianza degli array (cap.~\ref{ch:java}) \\
\hline
Python & nessun meccanismo dedicato: genericita' ``gratuita'' della tipizzazione dinamica, duck typing, \code{Protocol} come duck typing formalizzato (cap.~\ref{ch:python}) \\
\hline
C & nessuna vera genericita': macro (rischio di doppia valutazione), \code{void*}, \code{\_Generic} (overloading su un insieme chiuso, non generazione di codice) (cap.~\ref{ch:c}) \\
\hline
Vs. generics & i generics (Java/.NET) imitano solo il comportamento base dei template: mancano specializzazione, default argument, non-type/template-template argument (cap.~\ref{ch:confronto}) \\
\hline
Il filo che lega tutto & template, generics ed erasure, e duck typing sono tre risposte allo stesso problema, che differiscono solo per \emph{quando} verificano il ``duck test'' (cap.~\ref{ch:confronto}) \\
\hline
\end{tabular}
\end{center}

\section*{Fonti}
\addcontentsline{toc}{section}{Fonti}

\begin{itemize}
    \item Slide del corso: \emph{Templates in C++} (\code{20250922.TemplatesInCpp.pdf}), 40 pagine --- fonte primaria per i capitoli~\ref{ch:perche}--\ref{ch:specialization} e~\ref{ch:esercizio}.
    \item Wikipedia, \emph{Template (C++)}, \url{https://en.wikipedia.org/wiki/Template_(C++)} --- fonte per gli approfondimenti storici, le funzionalita' moderne del capitolo~\ref{ch:moderne} e l'apertura del confronto con i generics del capitolo~\ref{ch:confronto}.
    \item Wikipedia, \emph{Generics in Java}, \url{https://en.wikipedia.org/wiki/Generics_in_Java} --- fonte primaria del capitolo~\ref{ch:java}.
    \item Wikipedia, \emph{Generic programming}, \url{https://en.wikipedia.org/wiki/Generic_programming} --- fonte per la genericita' ``gratuita'' di Python (cap.~\ref{ch:python}) e per il C (cap.~\ref{ch:c}).
    \item Wikipedia, \emph{Duck typing}, \url{https://en.wikipedia.org/wiki/Duck_typing} --- fonte primaria del capitolo~\ref{ch:python}, in particolare del collegamento teorico template/duck typing.
    \item Bruce Eckel, \emph{Thinking in C++}, Vol.~2, cap.~5
          (\emph{Templates in Depth}), accessibile liberamente:
          \url{https://web.mit.edu/merolish/ticpp/TicV2.html\#_Toc53985706}
          --- fonte primaria del capitolo~\ref{ch:libro}.
    \item Cartella \code{prove/} (file \code{T01}--\code{T08},
          \code{es\_T01}--\code{es\_T02}): ogni esempio citato con
          ``Verificato in \dots'' e' stato \textbf{compilato ed eseguito}
          su questa macchina, non solo descritto.
\end{itemize}

\end{document}
